\subsection{正函数的叠积分}

在本节其余部分，如无相反的明确说明，我们将用 X 表示一个局部紧空间，用 \(\Lambda: t \mapsto \lambda_{t}\) 表示一个 \(\mu\) -适宜映射（T 到 \(\mathcal{M}_{+}(X)\) 中的），并用 \(\nu\) 表示 \(\Lambda\) 的积分。

命题 3. --- 设 \(f\) 是定义在 X 上的一个数值函数 \(\geqslant 0\) 。a) 下列不等式成立：

\[
\int^ {*} f (x) d \nu (x) \geqslant \int^ {\bullet} d \mu (t) \int^ {*} f (x) d \lambda_ {t} (x) \geqslant \int^ {\bullet} d \mu (t) \int^ {\bullet} f (x) d \lambda_ {t} (x). \tag {6}
\]

\begin{enumerate}
\def\labelenumi{\alph{enumi})}
\setcounter{enumi}{1}
\tightlist
\item
  若 \(\Lambda\) 是淡连续的，则
\end{enumerate}

\[
\int^ {*} f (x) d \nu (x) \geqslant \int^ {*} d \mu (t) \int^ {*} f (x) d \lambda_ {t} (x).\tag{7}
\]

\begin{enumerate}
\def\labelenumi{\alph{enumi})}
\setcounter{enumi}{2}
\tightlist
\item
  若 \(\lambda_t^\bullet (1) <   + \infty\) 局部 \(\mu\) -几乎处处成立，则
\end{enumerate}

\[
\int^ {\bullet} f (x) d \nu (x) \geqslant \int^ {\bullet} d \mu (t) \int^ {*} f (x) d \lambda_ {t} (x) = \int^ {\bullet} d \mu (t) \int^ {\bullet} f (x) d \lambda_ {t} (x). \tag {8}
\]

设 \(g\) 是 \(X\) 上的一个下半连续函数，满足 \(f \leqslant g\) 。对每个 \(t \in T\) ，

\[
\int^ {*} f (x) d \lambda_ {t} (x) \leqslant \int^ {*} g (x) d \lambda_ {t} (x),
\]

因此，由 (4) 及 §1 的命题 4，

\[
\int^ {\bullet} d \mu (t) \int^ {*} f (x) d \lambda_ {t} (x) \leqslant \int^ {\bullet} d \mu (t) \int^ {*} g (x) d \lambda_ {t} (x) = \int^ {*} g (x) d \nu (x).
\]

于是，不等式 (6) 中的第一个由 \(\int^{*}f(x)d\nu(x)\) 的定义（Ch. IV, §1, No.~3, 定义 3）得出，第二个则由它立即得出。若 \(\Lambda\) 是淡连续的，则不等式 (7) 用 (5) 代替 (4)，以类似的方式得到证明。

现在转向 (8) 的证明。映射 \(t \mapsto \lambda_t^*(1)\) 是可测的，且局部 \(\mu\) -几乎处处有限。因此，集 \(\mathfrak{K}\) （由 \(T\) 中使 \(t \mapsto \lambda_t^*(1)\) 到 \(K\) 上的限制为有限且连续的紧子集组成）是 \(\mu\) -稠密的，且 §2, No.~3 的命题 4 蕴涵存在正测度的一个可和族 \((\mu_\alpha)_{\alpha \in A}\) ，其支集属于 \(\mathfrak{K}\) ，使得 \(\mu = \sum_{\alpha \in A} \mu_\alpha\) 。映射 \(\Lambda\) 是 \(\mu_\alpha\) -适宜的（对每个 \(\alpha \in A\) ）；令 \(\nu_\alpha = \int \lambda_t d\mu_\alpha(t)\) 。命题 1 表明 \(\nu = \sum_{\alpha \in A} \nu_\alpha\) ，而把关系式 (4) 应用于测度 \(\mu_\alpha\) 和函数 1，就表明 \(\nu_\alpha\) 是有界测度（因为 \(\lambda_t^*(1)\) 在 \(\operatorname{Supp}(\mu_\alpha)\) 上有界）。于是对测度 \(\mu_\alpha\) 写出公式 (6) ，并将第一个成员中的符号 \(\int^*\) 换成 \(\int^\bullet\) ，由 §1 的命题 7，这是合法的；则

\[
\int^ {\bullet} f (x) d \nu_ {\alpha} (x) \geqslant \int^ {\bullet} d \mu_ {\alpha} (t) \int^ {*} f (x) d \lambda_ {t} (x) = \int^ {\bullet} d \mu_ {\alpha} (t) \int^ {\bullet} f (x) d \lambda_ {t} (x)
\]

（最后一个等式是由于 \(\lambda_t\) 局部几乎处处有界，以及 § 1 的命题 7）。然后对 \(\alpha\) 求和即得 (8) 中的不等式（§2, No.~2, 命题 1）。

如果不作任何与 c) 类似的假设，不等式 (8) 可能不成立（习题 2）。

推论 1. --- 设 \(f\) 是一个 \(\geqslant 0\) 的函数，定义在 \(X\) 上，并设 \(H\) 是由 \(t \in T\) （使 \(f\) 非 \(\lambda_t\) -可忽略者）组成的集。

\begin{enumerate}
\def\labelenumi{\alph{enumi})}
\item
  若 \(f \nu\) -可忽略，则 \(H\) 是局部 \(\mu\) -可忽略的。
\item
  若 \(f\) 是 \(\nu\) -可忽略的且 \(\Lambda\) 是淡连续的，则 \(H\) 是 \(\mu\) -可忽略的。
\item
  若 \(f\) 是局部 \(\nu\) -可忽略的且 \(\lambda_t^\bullet(1) < +\infty\) 局部 \(\mu\) -几乎处处成立，则 \(H\) 是局部 \(\mu\) -可忽略的。
\end{enumerate}

推论 2. --- 设 f 是定义在 X 上的一个函数 \(\geqslant0\) ，为 \(\nu\) -可测且 \(\nu\) -适度的。于是由 \(t\in T\) （使 f 非 \(\lambda_{t}\) -适度者）组成的集是局部 \(\mu\) -可忽略的（甚至是 \(\mu\) -可忽略的，如果 \(\Lambda\) 是淡连续的）。

因为，\(f\) 是一列函数 \(f_{n} \geqslant 0\) 之和，其中 \(f_{n}\) 在一个紧集 \(K_{n}\) 之外为零（对 \(n \geqslant 1\) ），而 \(f_{0}\) 是 \(\nu\) -可忽略的（§1, No.~2, 命题 6）；于是 \(f_{0}\) 是 \(\lambda_{t}\) -可忽略的，但 \(t\) 构成局部 \(\mu\) -可忽略集（甚至是 \(\mu\) -可忽略集，如果 \(\Lambda\) 是淡连续的）者除外（推论 1），从而命题立即得出。

命题 4. --- 设 f 是定义在 X 上、取值于拓扑空间 G 中的一个 \(\nu\) -可测函数，并设 M 是由 \(t \in T\) （使 f 非 \(\lambda_{t}\) -可测者）组成的集。

\begin{enumerate}
\def\labelenumi{\alph{enumi})}
\item
  假设 \(f\) 在一个 \(\nu\) -适度子集（\(X\) 中的）的余集上是常数；则 \(M\) 是局部 \(\mu\) -可忽略的。
\item
  假设 \(f\) 在一个 \(\nu\) -适度子集（\(X\) 中的）的余集上是常数，且 \(\Lambda\) 是淡连续的；则 \(M\) 是 \(\mu\) -可忽略的。
\item
  假设 \(\lambda_t^\bullet(1) < +\infty\) 局部 \(\mu\) -几乎处处成立；则 M 是局部 \(\mu\) -可忽略的。
\end{enumerate}

先证 a)（相应地，b)）。由于每个 \(\nu\) -可积集都含于一个 \(\nu\) -可积开集，函数 f 在可数个 \(\nu\) -可积开集之并的余集 B 上是常数。存在 X - B 的一个分割，由一个 \(\nu\) -可忽略集 N 与一列紧集 \((\mathrm{K}_{n})\) 组成，使得 f 到每个 \(K_{n}\) 上的限制连续。设 S 是由 \(t \in T\) （使 N 非 \(\lambda_{t}\) -可忽略者）组成的集：由命题 3 的推论 1，S 是局部 \(\mu\) -可忽略的（相应地，\(\mu\) -可忽略的）。集合 \(K_{n}, B, N\) 对 X 上每个测度都可测，且 f 到其中每一个上的限制是 \(\lambda_{t}\) -可测的（对每个 \(t \notin S\) ）。因此函数 f 是 \(\lambda_{t}\) -可测的（对每个 \(t \notin S\) ）（Ch. IV, §5, No.~10, 命题 16）。

为建立 c)，我们沿用命题 3 证明中的记号；由于 f 是 \(\nu\) -可测的，它对每个测度 \(\nu_{\alpha} \leqslant \nu\) 都可测。而这些测度是有界的，从而是适度的，于是由 a) 推出，M 是局部 \(\mu_{\alpha}\) -可忽略的（对每个 \(\alpha \in A\) ）。这蕴涵 M 是局部 \(\mu\) -可忽略的（§2, No.~2, 命题 1 的推论 2）。

命题 5. --- 设 f 是定义在 X 上的一个 \(\nu\) -可测正数值函数，并设 N 是由 \(t \in T\) （使 f 不同时为 \(\lambda_{t}\) -可测且 \(\lambda_{t}\) -适度者）组成的集。

\begin{enumerate}
\def\labelenumi{\alph{enumi})}
\tightlist
\item
  假设 \(f\) 是 \(\nu\) -适度的。于是集 \(\mathbf{N}\) 是局部 \(\mu\) -可忽略的，函数 \(t \mapsto \int^{\bullet} f(x) d\lambda_t(x)\) 是 \(\mu\) -可测的，且
\end{enumerate}

\[
\int^ {\bullet} f (x) d \nu (x) = \int^ {\bullet} d \mu (t) \int^ {\bullet} f (x) d \lambda_ {t} (x).\tag{9}
\]

\begin{enumerate}
\def\labelenumi{\alph{enumi})}
\setcounter{enumi}{1}
\tightlist
\item
  假设 \(f\) 是 \(\nu\) -适度的，且 \(\Lambda\) 是淡连续的。于是集 \(N\) 是 \(\mu\) -可忽略的，函数 \(t \mapsto \int^{*} f(x) d\lambda_{t}(x)\) 是 \(\mu\) -可测且 \(\mu\) -适度的，且
\end{enumerate}

\[
\int^ {*} f (x) d \nu (x) = \int^ {*} d \mu (t) \int^ {*} f (x) d \lambda_ {t} (x).\tag{10}
\]

\begin{enumerate}
\def\labelenumi{\alph{enumi})}
\setcounter{enumi}{2}
\tightlist
\item
  假设 \(\lambda_t^\bullet(1) < +\infty\) 局部 \(\mu\) -几乎处处成立。于是集 N 是局部 \(\mu\) -可忽略的，函数 \(t \mapsto \int^\bullet f(x) d\lambda_t(x)\) 是 \(\mu\) -可测的，且 (9) 成立。
\end{enumerate}

先在 \(f\) 是 \(\nu\) -适度的假设下证明 a)（相应地，b)）。关于集 \(N\) 的断言已经建立（命题 4，以及命题 3 的推论 2）。由 §1, No.~2 的命题 6，我们只需在下述每个特殊情形中证明 a)（相应地，b)）：

\begin{enumerate}
\def\labelenumi{\arabic{enumi})}
\item
  函数 \(f\) 是 \(\nu\) -可忽略的。
\item
  存在一个紧集 \(\mathbf{K}\) ，使得 \(f\) 在 \(\mathbf{K}\) 之外为零且 \(f\) 到 \(\mathbf{K}\) 上的限制连续。
\end{enumerate}

特殊情形 1) 已经处理过（命题 3 的推论 1）。为处理第二个情形，用 G 表示包含 K 的一个 \(\nu\) -可积开集，用 M 表示 f 的一个常数上界，用 h 表示下半连续函数 \(M\varphi_{G}\) ，并用 g 表示函数 h-f.~由于函数 f 在 X 上是上半连续的，g 是下半连续且正的。此外，f,g,h 都是 \(\nu\) -可积的。

于是把公式 (4)（相应地，(5)）应用于下半连续函数 h 与 g。由减法可见，函数

\[
t \mapsto \int^ {\bullet} f (x) d \lambda_ {t} (x) \quad (\text { resp. } \int^ {*} f (x) d \lambda_ {t} (x))
\]

是 \(\mu\) -可测的，且公式 (9)（相应地，(10)）成立。最后，在 b) 的假设下，函数 \(t \mapsto \int^{*} f(x) d\lambda_{t}(x)\) 具有有限的上积分，因而确实是 \(\mu\) -适度的。

为证 c)，我们沿用命题 3 证明中的测度 \(\mu_{\alpha}\) 与 \(\nu_{\alpha}\) ；由于 f 是 \(\nu_{\alpha}\) -可测且 \(\nu_{\alpha}\) -适度的，断言 a) 蕴涵 \(t \mapsto \int^{\bullet} f(x) d\lambda_{t}(x)\) 是 \(\mu_{\alpha}\) -可测的，且

\[
\int^ {\bullet} f (x) d \nu_ {\alpha} (x) = \int^ {\bullet} d \mu_ {\alpha} (t) \int^ {\bullet} f (x) d \lambda_ {t} (x).
\]

只需再对 \(\alpha\) 求和，并应用 §2, No.~2 的命题 1 与命题 2。

如果不假设 \(f\) 是 \(\nu\) -适度的，并且不作 c) 中的假设，则关系式 (9) 可能不成立（习题 3）。

推论. --- 设 f 是定义在 X 上、取值于 Banach 空间 F 或 \(\overline{R}\) 中的函数，且它是 \(\nu\) -可测并 \(\nu\) -适度的。为使 f 是 \(\nu\) -可积的，必须且只须

\[
\int^ {\bullet} d \mu (t) \int^ {\bullet} | \mathbf {f} (x) | d \lambda_ {t} (x) <   + \infty .
\]

这由命题 5 及可积性判别法（Ch. IV, §5, No.~6, 定理 5）立即得出。

\subsection{取值于 Banach 空间中的函数的叠积分}

定理 1. --- 设 f 是取值于 Banach 空间 F 或 \(\overline{R}\) 中的一个函数，并设 H 是由 \(t \in T\) （使 f 非 \(\lambda_{t}\) -可积者）组成的集。

\begin{enumerate}
\def\labelenumi{\alph{enumi})}
\tightlist
\item
  若 \(\mathbf{f}\) 是 \(\nu\) -可积的，则 \(\mathrm{H}\) 是局部 \(\mu\) -可忽略的，函数 \(t \mapsto \int \mathbf{f}(x)d\lambda_t(x)\) （对 \(t \notin \mathrm{H}\) 有定义）是本质 \(\mu\) -可积的，且
\end{enumerate}

\[
\int \mathbf {f} (x) d \nu (x) = \int d \mu (t) \int \mathbf {f} (x) d \lambda_ {t} (x).\tag{11}
\]

\begin{enumerate}
\def\labelenumi{\alph{enumi})}
\setcounter{enumi}{1}
\item
  若 \(\mathbf{f}\) 是 \(\nu\) -可积的且 \(\Lambda\) 是淡连续的，则此外 \(\mathbf{H}\) 还是 \(\mu\) -可忽略的，且函数 \(t \mapsto \int \mathbf{f}(x) d\lambda_t(x)\) （对 \(t \notin \mathbf{H}\) 有定义）是 \(\mu\) -可积的。
\item
  若 \(\lambda_t^\bullet(1) < +\infty\) 局部 \(\mu\) -几乎处处成立，则对本质 \(\nu\) -可积函数 f，a) 的结论仍成立。
\end{enumerate}

我们先来建立 a)（相应地，b)）。当 f 是正数值函数时，这一命题成立（命题 5）；若 f 是取值于 \(\overline{R}\) 的可积函数，则这一结果可以应用于正函数 \(f^{+}\) 与 \(f^{-}\) ，从而由减法立即推广到 f。剩下要处理取值于 F 的函数的情形。设 H 是 \(\mathcal{L}_{\mathrm{F}}^{1}(\nu)\) 中由 \(\mathcal{H}(\mathrm{X})\) 里的函数以 F 中元素为系数的线性组合所构成的子空间；关于实函数的结果立即蕴涵该命题对 H 的元素成立。而 H 在 \(\mathcal{L}_{\mathrm{F}}^{1}(\nu)\) 中稠密；因此，对每个 \(\mathbf{f} \in \mathcal{L}_{\mathrm{F}}^{1}(\nu)\) ，存在由 H 中元素组成的一个序列 \((\mathbf{f}_{n})\) ，它具有下列性质：

\begin{enumerate}
\def\labelenumi{\arabic{enumi})}
\item
  序列 \((\mathbf{f}_n)\) 平均收敛于 \(\mathbf{f}\) （在 \(\mathcal{L}_{\mathrm{F}}^{1}(\nu)\) 中），且 \(\nu\) -几乎处处收敛；
\item
  函数 \(g = |\mathbf{f}_0| + \sum_{n\in \mathbf{N}}|\mathbf{f}_{n + 1} - \mathbf{f}_n|\) 满足 \(\nu^{*}(g) < +\infty\) （Ch. IV, §3, No.~4, 定理 3）。
\end{enumerate}

设 \(N_{1}\) 是由 \(t \in T\) （使 \(\lambda_{t}^{*}(g) = +\infty\) 者）组成的集；由公式 (6)（相应地，(7)），\(N_{1}\) 是局部 \(\mu\) -可忽略的（相应地，\(\mu\) -可忽略的）。对 \(t \notin N_{1}\) ，诸 \(f_{n}\) 属于 \(\mathcal{L}_{F}^{1}(\lambda_{t})\) ，序列 \((f_{n})\) \(\lambda_{t}\) -几乎处处收敛，且对 \(\mathcal{L}_{F}^{1}(\lambda_{t})\) 中的平均收敛拓扑也收敛（Ch. IV, §3, No.~3, 命题 6）。设 M 是由 \(x \in X\) （使 \(f_{n}(x)\) 不收敛于 \(f(x)\) 者）组成的集：由于 M 是 \(\nu\) -可忽略的，集 \(N_{2}\) （由 \(t \in T\) （使 M 非 \(\lambda_{t}\) -可忽略者）组成）由命题 3 的推论 1 是局部 \(\mu\) -可忽略的（相应地，\(\mu\) -可忽略的）。

设 \(t\) 不属于 \(\mathrm{N}_1 \cup \mathrm{N}_2\) ；序列 \((\mathbf{f}_n)\) 在 \(\mathcal{L}_{\mathrm{F}}^{1}(\lambda_t)\) 中平均收敛，且 \(\lambda_t\) -几乎处处收敛于 \(\mathbf{f}\) 。因此 \(\mathbf{f} \in \mathcal{L}_{\mathrm{F}}^{1}(\lambda_t)\) 且 \(\int \mathbf{f} d\lambda_t = \lim_{n \to \infty} \int \mathbf{f}_n d\lambda_t\) （Ch. IV, §4, No.~1）。于是命题中的集 H 含于 \(N_{1} \cup N_{2}\) ；从而它是局部 \(\mu\) -可忽略的（相应地，\(\mu\) -可忽略的）。另一方面，函数 \(t \mapsto \int f d\lambda_{t}\) 局部 \(\mu\) -几乎处处等于一列 \(\mu\) -可测函数的极限；因此它是 \(\mu\) -可测的。最后，对每个 \(t \notin N_{1} \cup N_{2}\) 及每个 n，有

\[
\left| \int \mathbf {f} _ {n} (x) d \lambda_ {t} (x) \right| \leqslant \int^ {*} g (x) d \lambda_ {t} (x)
\]

这依据不等式 \(|\mathbf{f}_n| \leqslant g\) 及 Ch. IV, §4, No.~2 的命题 2。而函数 \(t \mapsto \int^{*} g(x) d\lambda_t(x)\) 由命题 5 是本质 \(\mu\) -可积的（相应地，\(\mu\) -可积的）。因此可以应用 Lebesgue 定理，它给出

\[
\int d \mu (t) \int \mathbf {f} (x) d \lambda_ {t} (x) = \lim _ {n \rightarrow \infty} \int d \mu (t) \int \mathbf {f} _ {n} (x) d \lambda_ {t} (x) = \lim _ {n \rightarrow \infty} \int \mathbf {f} _ {n} (x) d \nu (x).
\]

由于 \(\int \mathbf{f}_n(x) d\nu(x)\) 趋于 \(\int \mathbf{f}(x) d\nu(x)\) （当 \(n\) 趋于 \(\infty\) 时；依据对序列 \((\mathbf{f}_n)\) 所作的假设），关系式 (11) 得出，从而我们证明了 a)（相应地，b)）。

现在假设 \(\lambda_t^\bullet(1) < +\infty\) 局部 \(\mu\) -几乎处处成立，且 \(\mathbf{g}\) 是本质 \(\nu\) -可积函数。设 \(\mathbf{f}\) 是一个 \(\nu\) -可积函数，使得 \(\mathbf{g} = \mathbf{f}\) 局部 \(\nu\) -几乎处处成立（§1, No.~3）。于是 \(\mathbf{g} = \mathbf{f}\) 对 \(\lambda_t\) 几乎处处成立，除去由 \(t\) 组成的一个局部 \(\mu\) -可忽略集 \(\mathrm{P}\) 的情形（命题 3 的推论 1 c)）。因此 \(\int \mathbf{g} d\lambda_t = \int \mathbf{f} d\lambda_t\) 对所有 \(t \notin \mathrm{P} \cup \mathrm{H}\) 成立，这就完成了证明。

注. --- 设 \(\Lambda: t \mapsto \lambda_{t}\) 是一个 \(\mu\) -适宜映射（T 到 \(\mathcal{M}_{+}(X)\) 中的）。如果一个映射 \(\Lambda': t \mapsto \lambda_{t}'\) （T 到 \(\mathcal{M}_{+}(X)\) 中的）与 \(\Lambda\) 局部 \(\mu\) -几乎处处相等，则由定义立即可见 \(\Lambda'\) 也是 \(\mu\) -适宜的，且 \(\Lambda\) 与 \(\Lambda'\) 有相同的积分。现在设 \(H: t \mapsto \eta_{t}\) 是一个取值于 \(\mathcal{M}_{+}(X)\) 中的函数，它只在局部 \(\mu\) -几乎处处有定义；我们仍称 H 是 \(\mu\) -适宜的，如果它局部 \(\mu\) -几乎处处等于一个映射 \(\Lambda: t \mapsto \lambda_{t}\) ，后者处处有定义且 \(\mu\) -适宜。这时我们置 \(\int \eta_{t} d\mu(t) = \int \lambda_{t} d\mu(t)\) ，这一定义不依赖于所用的函数 \(\Lambda\) 。前面各 No. 中证明的命题均可推广到局部 \(\mu\) -几乎处处有定义的 \(\mu\) -适宜函数，验证工作留给读者。

\subsection{普遍可测函数}

定义 2. --- T 到拓扑空间 F 中的映射 f 称为普遍可测的，如果它对 T 上每个正测度 \(\mu\) 都是 \(\mu\) -可测的。

T 的子集中，其特征函数为普遍可测者，称为普遍可测集。它们构成 T 上的一个 σ-集环（tribe）（Ch. IV, §5, No.~4, 定理 2 的推论 2），该 σ-集环包含 Borel 集（同处引证，推论 3），并且当 T 可度量化时还包含 Souslin 集（Ch. IV, §5, No.~1, 命题 3 的推论 2）。T 到可度量化且可分的拓扑空间 F 中的映射 f 为普遍可测的，必须且只须 F 中每个闭球在 f 下的逆像是 T 的普遍可测子集（Ch. IV, §5, No.~5, 定理 4）。

命题 6. --- T 到拓扑空间 F 中的映射 f 为普遍可测的，必须且只须 f 对 T 上每个具有紧支集的正测度都可测。

该条件显然是必要的；另一方面它也是充分的，因为每个正测度 \(\mu\) 都是一族具有紧支集的测度之和（ \(\S2\) ，No.~3，命题 4）：于是结论由 \(\S2\) ，No.~2 的命题 2 得出。

命题 7. --- 设 \(\mu\) 是 T 上的一个正测度，\(f\) 是 T 到拓扑空间 F 中的一个 \(\mu\) -可测映射。则存在 T 到 F 中的一个普遍可测映射 \(f'\) ，使得 \(f = f'\) 局部 \(\mu\) -几乎处处成立。

设 \(\mathfrak{K}\) 是 T 中使 \(f\) 在 K 上的限制为连续的紧集组成的集；由于 \(\mathfrak{K}\) 是 \(\mu\) -稠密的（Ch. IV, §5, No.~10, 命题 15），存在一个局部可数族 \((\mathrm{K}_i)_{i\in I}\) （由 \(\mathfrak{K}\) 中两两不相交的元素组成），使得集 \(\mathrm{N} = \mathrm{T} - \bigcup_{i\in I}\mathrm{K}_i\) 是局部 \(\mu\) -可忽略的（Ch. IV, §5, No.~9,

命题 14）。设 \(x\) 是 \(\mathbf{F}\) 的一个元素；置

\[
f ^ {\prime} (t) = f (t) \quad \text { if } t \in \bigcup_ {i \in I} K _ {i},
\]

\[
f ^ {\prime} (t) = x \quad \text { if } t \in \mathbb {N}.
\]

函数 \(f\) 与 \(f'\) 局部 \(\mu\) -几乎处处相等。另一方面，\(\mathrm{N} \cap \mathrm{K}\) 是 \(\mathrm{K}\) 的一个 Borel 子集，其中 \(\mathrm{K}\) 是 \(\mathrm{T}\) 中任意的紧集，因为族 \((\mathrm{K}_i)\) 是局部可数的。由此得出 \(\mathrm{N}\) 是普遍可测集，且 \(f'\) 是普遍可测函数（Ch. IV, §5, No.~10, 命题 16）。

\subsection{扩散}

定义 3. --- 设 X 是一个局部紧空间，\(\Lambda : t \mapsto \lambda_t\) 是 T 到 \(\mathcal{M}_{+}(X)\) 中的一个映射。称映射 \(\Lambda\) 是 T 在 X 中的一个扩散，如果 \(\Lambda\) 对 T 上每个具有紧支集的正测度都是适宜的。称扩散 \(\Lambda\) 是有界的，如果所有测度 \(\lambda_{t}\) 都有界且 \(\sup_{t\in T}\|\lambda_{t}\|<+\infty\) ；这个量称为 \(\Lambda\) 的范数，记作 \(\|\Lambda\|\) 。

下面的命题只是定义的转述：

命题 8. --- 映射 \(\Lambda: t \mapsto \lambda_t\) （T 到 \(\mathcal{M}_+(\mathrm{X})\) 中的）是一个扩散，必须且只须下列条件成立：

\begin{enumerate}
\def\labelenumi{\arabic{enumi})}
\item
  对每个下半连续函数 \(f \geqslant 0\) （定义在 \(X\) 上），函数 \(t \mapsto \lambda_t^\bullet(f)\) 在 \(T\) 上是普遍可测的。
\item
  对每个函数 \(g \in \mathcal{K}_{+}(X)\) ，函数 \(t \mapsto \lambda_t(g)\) 在 \(T\) 中是局部有界的。
\item
  对每个定义在 X 上的下半连续函数 \(f \geqslant 0\) 以及 T 上每个具有紧支集的正测度 \(\mu\) ，下列关系式成立，其中 \(\nu\) 表示 \(\int \lambda_{t} d\mu(t)\) ：
\end{enumerate}

\[
\int^ {\bullet} f (x) d \nu (x) = \int^ {\bullet} d \mu (t) \int^ {\bullet} f (x) d \lambda_ {t} (x).\tag{12}
\]

假设 \(\Lambda\) 是一个扩散。则由适宜映射的定义（No.~1，定义 1）及命题 6，条件 1) 成立；由公式 (4)，条件 3) 成立，因为 \(\Lambda\) 是 \(\mu\) -适宜的。设 \(g \in \mathcal{K}_{+}(X)\) ，\(u\) 是函数 \(t \mapsto \lambda_t(g)\) （由 1)，它是普遍可测的）；假设 \(u\) 不是局部有界的。那么将存在一个紧集 \(K\) 使 \(u\) 在 \(K\) 上不有界，从而将存在元素的一个序列 \((t_n)\) （取自 \(K\) ），使 \(u(t_n) \geqslant n^2\) 对所有 \(n \geqslant 1\) 成立；于是 \(u\) 对具有紧支集的测度 \(\mu = \sum_{n \geqslant 1} \frac{1}{n^2} \varepsilon_{t_n}\) 不是可积的，这与关于 \(\Lambda\) 的假设相矛盾，该假设蕴涵 \(t \mapsto \lambda_t(g)\) 对每个具有紧支集的正测度可积。因此上述三个条件是必要的。反之，条件 1) 与 2) 蕴涵 \(\Lambda\) 对每个具有紧支集的测度 \(\mu\) 都是标量本质 \(\mu\) -可积的。由于每个测度 \(\mu' \geqslant 0\) （被一个具有紧支集的测度 \(\mu\) 优超者）本身也具有紧支集，条件 1) 与 3) 表明 \(\Lambda\) 对每个具有紧支集的正测度都是 \(\mu\) -适宜的，这正是所要的结果。

命题 9. --- 设 \(\Lambda : t \mapsto \lambda_t\) 是 T 到 \(\mathcal{M}_+(\mathrm{X})\) 中的一个映射，使得函数 \(t \mapsto \lambda_t(g)\) 对每个 \(g \in \mathcal{K}_+(\mathrm{X})\) 都在 T 上普遍可测且局部有界。在下列每一种情形，都可以断言 \(\Lambda\) 是一个扩散：

\begin{enumerate}
\def\labelenumi{\alph{enumi})}
\item
  X 的拓扑具有可数基；
\item
  \(\Lambda\) 关于淡拓扑是普遍可测的。
\end{enumerate}

因为，设 \(\mu\) 是 T 上一个具有紧支集的正测度；映射 \(\Lambda\) 是标量本质 \(\mu\) -可积的，因而只要 a) 或 b) 之一成立，它就是 \(\mu\) -适宜的（No.~1，命题 2）。

在本节其余部分，我们将采用以下记号：我们将用 \(\langle\eta,h\rangle\) 表示正测度 \(\eta\) 下的正 \(\eta\) -可测函数 h 的本质上积分。映射 \(\Lambda:t\mapsto\lambda_{t}\) 将是 T 在 X 中的一个扩散。如果 f 是定义在 X 上的一个正的普遍可测函数，我们将用 \(\Lambda f\) 表示映射 \(t\mapsto\lambda_{t}^{\bullet}(f)\) 。如果 \(\mu\) 是 T 上使 \(\Lambda\) 为标量本质 \(\mu\) -可积的正测度，我们将用 \(\mu\Lambda\) 表示测度 \(\int\lambda_{t}d\mu(t)\) 。于是积分的定义取如下形式

\[
\langle \mu \Lambda , f \rangle = \langle \mu , \Lambda f \rangle \quad \text { for } f \in \mathscr {K} _ {+} (\mathrm{X}).
\]

我们将说 T 上的正测度 \(\mu\) 属于 \(\Lambda\) 的定义域，如果 \(\Lambda\) 是 \(\mu\) -适宜的：这等价于说（依据命题 8）\(\Lambda\) 是标量本质 \(\mu\) -可积的，且 \(\langle\mu'\Lambda,f\rangle=\langle\mu',\Lambda f\rangle\) 对每个正测度 \(\mu'\leqslant\mu\) 及每个正的下半连续函数 f 成立。

命题 10. --- 设 f, g 是 X 上两个正的普遍可测函数，a 是一个数 \(\geqslant 0\) ，\(\mu\) 和 \(\nu\) 是 T 上两个正测度。那么：

\begin{enumerate}
\def\labelenumi{\alph{enumi})}
\item
  \(\Lambda (f + g) = \Lambda f + \Lambda g, \Lambda (af) = a\Lambda f.\)
\item
  如果 \(\mu\) 和 \(\nu\) 属于 \(\Lambda\) 的定义域，那么 \(\mu + \nu\) 和 \(a\mu\) 也属于其定义域，并且有 \((\mu + \nu)\Lambda = \mu \Lambda + \nu \Lambda\) ，\((a\mu)\Lambda = a(\mu \Lambda)\) 。
\end{enumerate}

唯一不明显之点是 \(\mu + \nu\) 属于 \(\Lambda\) 的定义域，这一点可以这样处理：注意到每个被 \(\mu + \nu\) 优超的正测度都具有 \(\mu' + \nu'\) 的形式，其中 \(\mu' \leqslant \mu\) ，\(\nu' \leqslant \nu\) （“分解引理”，Ch. II, §1, No.~1）。另见下一命题。

命题 11. --- 正测度 \(\mu\) （\(T\) 上的）属于 \(\Lambda\) 的定义域，必须且只须 \(\Lambda\) 是标量本质 \(\mu\) -可积的。

该条件显然是必要的。反之，设它成立，并设 f 是定义在 X 上的一个正的下半连续函数。函数 \(\Lambda f\) 是普遍可测的，因而是 \(\mu\) -可测的。我们将证明 \(\langle\mu,\Lambda f\rangle=\langle\mu\Lambda,f\rangle\) ；由于这个等式对每个正测度 \(\mu'\leqslant\mu\) 也成立（因为 \(\Lambda\) 也是标量本质 \(\mu'\) -可积的），便可推出 \(\Lambda\) 是 \(\mu\) -适宜的。

设 \((\mu_i)_{i\in I}\) 是一个由具有紧支集的正测度组成的可和族，使得 \(\mu = \sum_{i\in I}\mu_i\) （§2, No.~3, 命题 4）；于是测度族 \(\mu_{i}\Lambda\) 是可和的，且 \(\mu \Lambda = \sum_{i\in I}\mu_{i}\Lambda\) （No.~1, 命题 1 的推论）。从而 \(\langle \mu \Lambda ,f\rangle = \sum_{i\in I}\langle \mu_i\Lambda ,f\rangle\) （§2, No.~2, 命题 1）；但 \(\Lambda\) 是 \(\mu_{i}\) -适宜的，故 \(\langle \mu_i\Lambda ,f\rangle = \langle \mu_i,\Lambda f\rangle\) 。再应用 §2 的命题 1，就得到所要的等式：

\[
\langle \mu \Lambda , f \rangle = \sum_ {i \in I} \langle \mu_ {i} \Lambda , f \rangle = \sum_ {i \in I} \langle \mu_ {i}, \Lambda f \rangle = \langle \mu , \Lambda f \rangle .
\]

推论 1. --- 如果 \(\Lambda\) 是一个有界扩散，那么每个有界正测度 \(\mu\) 都属于 \(\Lambda\) 的定义域，且 \(\|\mu\Lambda\| \leqslant \|\mu\|\|\Lambda\|\) 。

推论 2. --- 假设 \(\mu\) 是一个可和族 \((\mu_{\alpha})_{\alpha \in \mathbb{A}}\) 之和，该族由属于 \(\Lambda\) 的定义域的正测度组成。\(\mu\) 属于 \(\Lambda\) 的定义域，必须且只须测度族 \(\mu_{\alpha}\Lambda\) 可和，此时 \(\mu \Lambda = \sum \mu_{\alpha}\Lambda\) 。

只需应用 No.~1 的命题 1 的推论。

命题 5 用扩散的语言表述，就取如下的形式：

命题 12. --- 设 \(\mu\) 是 T 上属于 \(\Lambda\) 的定义域的一个正测度，\(f\) 是定义在 X 上的一个普遍可测函数 \(\geqslant 0\) 。如果 \(f\) 关于测度 \(\mu\Lambda\) 是适度的，或者如果诸测度 \(\lambda_t\) 有界，那么函数 \(\Lambda f\) 是 \(\mu\) -可测的，并且

\[
\langle \mu \Lambda , f \rangle = \langle \mu , \Lambda f \rangle .\tag{13}
\]

推论. --- 如果 X 在无穷远处可数，或者如果诸测度 \(\lambda_{t}\) 有界，那么函数 \(\Lambda f\) 对每个定义在 X 上的普遍可测函数 \(f \geqslant 0\) 都在 T 上普遍可测，且 (13) 成立。

\subsection{有界扩散的复合}

命题 13. --- 设 T, X, Y 是三个局部紧空间，\(\Lambda : t \mapsto \lambda_t\) 是 T 在 X 中的一个有界扩散，H: \(x \mapsto \eta_x\) 是 X 在 Y 中的一个有界扩散。则映射 \(t \mapsto \lambda_t\mathrm{H}\) 是 T 在 Y 中的一个有界扩散，记作 \(\Lambda\mathrm{H}\) ，并且

\[
\| \Lambda H \| \leqslant \| \Lambda \| \| H \|.\tag{14}
\]

设 \(f\) 是一个普遍可测函数 \(\geqslant 0\) （定义在 \(Y\) 上），\(\mu\) 是 \(T\) 上一个测度。假设 \(\mu\) 属于 \(\Lambda\) 的定义域，且 \(\mu \Lambda\) 属于 H 的定义域；则 \(\mu\) 属于 \(\Lambda \mathrm{H}\) 的定义域，并且

\[
\begin{array}{l} \langle \mu (\Lambda \mathrm{H}), f \rangle = \langle \mu \Lambda , \mathrm{H} f \rangle = \langle \mu , \Lambda \mathrm{H} f \rangle ; \\ (\mu \Lambda) \mathrm{H} = \mu (\Lambda \mathrm{H}); \quad \Lambda (\mathrm{H} f) = (\Lambda \mathrm{H}) f. \end{array}\tag{15}
\]

置 \(\gamma_{t} = \lambda_{t}H\) ；我们将用 \(\Gamma\) 表示映射 \(\Lambda H\) （T 到 \(\mathcal{M}_{+}(Y)\) 中的），用 \(\Gamma f\) 表示函数 \(t \mapsto \langle \gamma_{t}, f \rangle\) （这是记号的滥用，因为我们尚不知道 \(\Gamma\) 是否为扩散）。由 (13)，\(\langle \gamma_{t}, f \rangle = \langle \lambda_{t}H, f \rangle = \langle \lambda_{t}, Hf \rangle\) ；由于函数 Hf 在 X 上是正的且普遍可测的（命题 12 的推论），首先得出 \(\Gamma f = \Lambda(Hf)\) ，继而得出 \(\Gamma f\) 在 T 上普遍可测（同处引证）。显然，所有测度 \(\gamma_{t}\) 的总质量至多等于 \(\|\Lambda\| \|H\|\) 。因此 \(\Gamma g\) 对每个函数 \(g \in \mathcal{K}_{+}(Y)\) 都是普遍可测且有界的；从而 \(\Gamma\) 对 T 上每个有界测度、特别地对每个具有紧支集的测度，是标量本质可积的。更一般地，如果 \(\mu\) 是 \(\Lambda\) 的定义域中的一个测度，使得 \(\mu\Lambda\) 属于 H 的定义域，那么对 \(g \in \mathcal{K}_{+}(Y)\) ，

\[
\langle \mu , \Gamma g \rangle = \langle \mu , \Lambda (\mathrm{H} g) \rangle = \langle \mu \Lambda , \mathrm{H} g \rangle = \langle (\mu \Lambda) \mathrm{H}, g \rangle .
\]

由于最后的量是有限的，我们看出 \(\Gamma\) 是标量本质 \(\mu\) -可积的。让我们用 \(\mu\Gamma\) 表示积分 \(\int\gamma_{t}d\mu(t)\) （这是记号的滥用，因为我们尚不知道 \(\Gamma\) 是否为扩散）。于是前面的关系式可以写成

\[
\langle \mu \Gamma , g \rangle = \langle (\mu \Lambda) \mathrm{H}, g \rangle ,
\]

或者写成 \(\mu \Gamma = (\mu \Lambda)\mathrm{H}\) ，因为 \(g\) 在 \(\mathcal{K}_+(\mathrm{Y})\) 中是任意的。

再次考虑普遍可测函数 \(f \geqslant 0\) 。我们有

\[
\langle \mu \Gamma , f \rangle = \langle (\mu \Lambda) \mathrm{H}, f \rangle = \langle \mu \Lambda , \mathrm{H} f \rangle = \langle \mu , \Lambda (\mathrm{H} f) \rangle = \langle \mu , \Gamma f \rangle .
\]

当 f 下半连续且 \(\mu\) 跑遍具有紧支集的正测度所成的集时，这些关系式表明 \(\Gamma\) 是 T 在 Y 中的一个扩散。于是该论断不过是把上面证明过程中得到的关系式明显化而已。

定义 4. --- 记号沿用命题 13，扩散 \(\Lambda H\) 称为有界扩散 \(H\) 与 \(\Lambda\) 的复合扩散（或复合）。

设 \(X_{1}, X_{2}, X_{3}, X_{4}\) 是四个局部紧空间，\(\Lambda_{1}, \Lambda_{2}, \Lambda_{3}\) 是三个有界扩散，分别为 \(X_{1}\) 在 \(X_{2}\) 中、\(X_{2}\) 在 \(X_{3}\) 中、\(X_{3}\) 在 \(X_{4}\) 中的扩散。由命题 13 立即得出

\[
(\Lambda_ {1} \Lambda_ {2}) \Lambda_ {3} = \Lambda_ {1} (\Lambda_ {2} \Lambda_ {3}).
\]

因此，对于扩散的复合，我们将使用这些不带括号的记号。

例子. --- 设 u 是 T 到 X 中的一个普遍可测映射，v 是 X 到 Y 中的一个普遍可测映射；由命题 2 b)，用下列公式定义扩散 \(\Lambda\) 和 H：

\[
\lambda_ {t} = \varepsilon_ {u (t)}, \eta_ {x} = \varepsilon_ {v (x)};
\]

则扩散 \(\Gamma = \Lambda H\) 由下式给出：

\[
\gamma_ {t} = \varepsilon_ {(v \circ u) (t)}.
\]

因此要小心注意：扩散的复合次序与通常函数复合的次序相反。

\section{正点测度的积分}

\subsection{点测度族}

设 X 和 T 是两个局部紧空间，\(\pi\) 是 T 到 X 中的一个映射，\(g\) 是定义在 T 上的一个有限数值函数 \(\geqslant 0\) ；这两个函数定义了映射 \(t \mapsto \lambda_t = g(t)\varepsilon_{\pi(t)}\) （T 到 X 上测度所成的空间 \(\mathcal{M}(X)\) 中的），使得对每个 \(t \in T\) ，\(\lambda_t\) 或者是一个点测度（Ch. III, §2, No.~4），或者等于 0。如果 \(f\) 是定义在 X 上的一个数值函数 \(\geqslant 0\) ，那么 \(\int^{*} f(x) d\lambda_t(x) = \int^{\bullet} f(x) d\lambda_t(x) = f(\pi(t))g(t)\) （回忆我们的约定：当 \(g(t) = 0\) 且 \(f(\pi(t)) = +\infty\) 时取这个乘积为 0）。定义在 X 上的每个函数（取值于拓扑空间）都是 \(\lambda_t\) -可测的（对每个 \(t \in T\) ）。X 到 Banach 空间 F 中的每个映射 \(\mathbf{f}\) 都是 \(\lambda_t\) -可积的（对所有 \(t \in T\) ），且 \(\int \mathbf{f}(x) d\lambda_t(x) = \mathbf{f}(\pi(t))g(t)\) 。最后，如果 \(f\) 是定义在 X 上的任意数值函数，那么 \(f\) 为 \(\lambda_t\) -可积必须且只须 \(f(\pi(t))g(t)\) 有限，此时 \(\int f(x) d\lambda_t(x) = f(\pi(t))g(t)\) 。

定义 1. --- 设 \(\mu\) 是 T 上一个正测度。称对 \((\pi, g)\) 为 \(\mu\) -适应对，如果下列条件成立：

\(1^{\circ}\) 函数 \(\pi\) 和 \(g\) 都是 \(\mu\) -可测的。

\(2^{\circ}\) 对每个函数 \(f \in \mathcal{K}(X)\) ，映射 \(t \mapsto f(\pi(t))g(t)\) 都是本质 \(\mu\) -可积的。

命题 1. --- 如果对 \((\pi, g)\) 是 \(\mu\) -适应的，那么映射 \(\Lambda: t \mapsto \lambda_t = g(t) \varepsilon_{\pi(t)}\) （T 到 \(\mathcal{M}_+(\mathrm{X})\) 中的）是标量本质 \(\mu\) -可积的、淡 \(\mu\) -可测的且 \(\mu\) -适宜的。反之，如果 \(\Lambda\) 是标量本质 \(\mu\) -可积且淡 \(\mu\) -可测的，那么函数 g 是 \(\mu\) -可测的，且 \(\pi\) 在集 S（由 \(t \in T\) 中使 \(g(t) \neq 0\) 者组成的）上的限制是 \(\mu\) -可测的。

假设对 \((\pi,g)\) 是 \(\mu\) -适应的；对每个函数 \(f \in \mathcal{K}(X)\) ，函数 \(t \mapsto \langle f, \lambda_t \rangle = f(\pi(t))g(t)\) 都是本质 \(\mu\) -可积的。让我们证明 \(t \mapsto \lambda_t\) 是淡 \(\mu\) -可测的。因为，我们首先注意到，如果 \(\pi\) 和 g 连续，那么映射 \(t \mapsto \lambda_t\) 是淡连续的。在一般情形，T 中使 \(\pi\) 和 g 在 K 上的限制都连续的紧子集 K 所成的集是 \(\mu\) -稠密的（Ch. IV, §5, No.~10, 命题 15）；如果 K 是这样一个集，那么 \(t \mapsto \lambda_t\) 在 K 上的限制是淡连续的，由此得出陈述的第一个断言。§3, No.~1 的命题 2 表明 \(\Lambda\) 是 \(\mu\) -适宜的。

反之，假设 \(\Lambda\) 是标量本质 \(\mu\) -可积且淡 \(\mu\) -可测的；那么它是 \(\mu\) -适宜的（§3, No.~1, 命题 2）。由于函数 1 在 X 上是下半连续的，函数 \(t \mapsto \lambda_t(1) = g(t)\) 是 \(\mu\) -可测的（§3, 定义 1）。因此集 S 是可测的（Ch. IV, §5, No.~5, 命题 7）。紧集 K ⊂ S 所成的集 \(\mathfrak{K}\) （这些 K 使 \(g|K\) 连续且 \(\Lambda|K\) 淡连续）在 S 中是 \(\mu\) -稠密的（Ch. IV, §5, No.~10, 命题 15）；如果 K ∈ \(\mathfrak{K}\) ，那么映射 \(t \mapsto \varepsilon_{\pi(t)} = \frac{1}{g(t)} \lambda_t\) 在 K 上的限制是淡连续的，这就蕴涵 \(\pi|K\) 的连续性（Ch. III, §1, No.~9, 命题 13）。由于 \(\mathfrak{K}\) 在 S 中 \(\mu\) -稠密，\(\pi\) 在 S 上的限制是 \(\mu\) -可测的。

我们将利用下面的引理：

引理. --- 设 T 和 X 是两个拓扑空间，\(\pi\) 是 T 到 X 中的一个逆紧连续映射（GT, I, §10, No.~1, 定义 1）。设 g 是定义在 T 上的一个下半连续数值函数。对每个 \(x \in X\) ，令 \(f(x)\) 为函数 \(g(t)\) 在集 \(\overline{\pi}^{-1}(x)\) 中的下确界（下确界等于 \(+\infty\) ，如果 \(\overline{\pi}^{-1}(x) = \varnothing\) ；参见 S, III, §1, No.~9）。则 f 在 X 上是下半连续的。

对每个（有限）实数 \(a\) ，用 \(\mathrm{B}_a\) 表示由 \(x \in \mathrm{X}\) 中使 \(f(x) \leqslant a\) 者组成的集，用 \(\mathrm{A}_a\) 表示由 \(t \in \mathrm{T}\) 中使 \(g(t) \leqslant a\) 者组成的集；问题归结为证明 \(\mathrm{B}_a\) 是闭的（GT, IV, §6, No.~2, 命题 1）。现在，\(\mathrm{A}_a\) 是闭的（同处引证），且逆紧映射 \(\pi\) 是闭的（GT, I, §10, No.~1, 命题 1）；于是我们被归结为证明 \(\pi(\mathrm{A}_a) = \mathrm{B}_a\) 。明显的关系 \(f(\pi(t)) \leqslant g(t)\) （对所有 \(t \in \mathrm{T}\) ）蕴涵 \(\pi(\mathrm{A}_a) \subset \mathrm{B}_a\) 。另一方面，设 \(x \in \mathrm{B}_a\) ；集 \(\overline{\pi}^{-1}(x)\) 是拟紧的（GT, I, §10, No.~2, 定理 1）且非空，因此存在 \(t \in \overline{\pi}^{-1}(x)\) 使 \(g(t) = \inf_{u\in \overline{\pi}^{-1} (x)}g(u) = f(x)\) （GT, IV, §6, No.~2, 定理 3）；从而 \(t\in \mathbf{A}_a\) 且 \(\pi (t) = x\)。

\subsection{关于点测度积分的正函数上积分}

我们将看到，当 \((\pi,g)\) 是一个 \(\mu\) -适应对时，把 §3 的诸命题应用于族 \(t\mapsto\lambda_{t}=g(t)\varepsilon_{\pi(t)}\) （由命题 1，它是 \(\mu\) -适宜的）所得到的结果可以被加强。

定理 1. --- 设 \((\pi, g)\) 是一个 \(\mu\) -适应对，并令

\[
\nu = \int g (t) \varepsilon_ {\pi (t)} d \mu (t).
\]

对每个数值函数 \(f \geqslant 0\) （定义在 \(X\) 上），

\[
\int^ {\bullet} f (x) d \nu (x) = \int^ {\bullet} f (\pi (t)) g (t) d \mu (t).\tag{1}
\]

\begin{enumerate}
\def\labelenumi{\Alph{enumi})}
\tightlist
\item
  首先假设测度 \(\mu\) 具有紧支集 \(K\) ，且限制到 \(K\) 上时函数 \(g\) 和 \(\pi\) 都连续。由 §3, No.~1 的公式 (4)，\(\nu^{\bullet}(1) = \int_{K} g(t) d\mu(t) < +\infty\) ，故出现在公式 (1) 中的所有测度都是有界的。因此我们可以在第一成员中把 \(\int^{\bullet}\) 换成 \(\int^{*}\) 。鉴于 §3, No.~2 的公式 (6)，问题归结为证明
\end{enumerate}

\[
\int^ {*} f (x) d \nu (x) \leqslant \int^ {\bullet} f (\pi (t)) g (t) d \mu (t),\tag{2}
\]

其中第二成员中的记号 \(\int^{\bullet}\) 又可以换成 \(\int^{*}\) 。由上积分的定义，只需验证不等式

\[
\int^ {*} f (x) d \nu (x) \leqslant \int^ {*} h (t) d \mu (t)\tag{3}
\]

对每个下半连续函数 \(h\) （\(\mathrm{T}\) 上的）——它 \(\geqslant\) 函数 \(t \mapsto f(\pi(t))g(t)\) 。现在，设 \(\varepsilon\) 是一个数 \(> 0\) ，\(u\) 是函数 \((h + \varepsilon)/g\) ，它在 \(\mathrm{K}\) 中是下半连续的。如果 \(t \in \overline{\pi}^{-1}(\{x\}) \cap \mathrm{K}\) ，那么 \(u(t) \geqslant f(x)\) ：当 \(g(t) = 0\) 时这是明显的，因为那时 \(u(t) = +\infty\) ；如果 \(g(t) > 0\) ，那么

\[
u (t) g (t) = h (t) + \varepsilon \geqslant f (\pi (t)) g (t) = f (x) g (t),
\]

由此得出所要的不等式。在这些条件下，对每个 \(x \in X\) ，令 \(v(x)\) 为 \(u(t)\) 对 \(t \in \overline{\pi}^{-1}(\{x\}) \cap K\) 的下确界。由前所述，函数 v \(\geqslant f\) ；由引理（应用于 \(\pi\) 在 K 上的限制），它在 X 上是下半连续的；且 \(v(\pi(t))g(t) \leqslant h(t) + \varepsilon\) 对所有 \(t \in K\) 成立（回忆：按约定，若 \(g(t) = 0\) ，则第一成员为零）。于是对 v 应用 §3, No.~1 的公式 (4)。我们得到：

\[
\begin{array}{l} \int^ {*} f (x) d \nu (x) \leqslant \int^ {*} v (x) d \nu (x) \\ = \int^ {*} v (\pi (t)) g (t) d \mu (t) \leqslant \int_ {K} ^ {*} (h (t) + \varepsilon) d \mu (t) \\ = \int^ {*} h (t) d \mu (t) + \varepsilon \mu (1). \end{array}\tag{4}
\]

由于测度 \(\mu\) 有界且 \(\varepsilon\) 是任意的，就得出不等式 (3)。

\begin{enumerate}
\def\labelenumi{\Alph{enumi})}
\setcounter{enumi}{1}
\tightlist
\item
  现在转到一般情形。由于映射 \(t \mapsto (\pi(t), g(t))\) （\(T\) 到 \(X \times R_{+}\) 中的）是 \(\mu\) -可测的（Ch. IV, §5, No.~3, 定理 1），集 \(\mathfrak{K}\) （由紧子集 \(K\) （\(T\) 中的）组成，使 \(\pi\) 和 \(g\) 到 \(K\) 上的限制都连续）是 \(\mu\) -稠密的（Ch. IV, §5, No.~10, 命题 15）。由 §2, No.~3 的命题 4，\(\mu\) 是一个可和族 \((\mu_{\alpha})_{\alpha \in A}\) 之和，这些测度的支集都是 \(\mathfrak{K}\) 的元素；对 \((\pi, g)\) 是 \(\mu_{\alpha}\) -适应的（对每个 \(\alpha \in A\) ），令 \(\nu_{\alpha}\) 为测度 \(\int g(t)\varepsilon_{\pi(t)} d\mu_{\alpha}(t)\) 。于是由 A)，
\end{enumerate}

\[
\int^ {\bullet} f (x) d \nu_ {\alpha} (x) = \int^ {\bullet} f (\pi (t)) g (t) d \mu_ {\alpha} (t).\tag{5}
\]

但诸 \(\nu_{\alpha}\) 构成一个可和族，其和等于 \(\nu\) （§3, No.~1, 命题 1 的推论）。因此，由 §2, No.~2 的命题 1，

\[
\int^ {\bullet} f (x) d \nu (x) = \sum_ {\alpha \in \mathrm{A}} \int^ {\bullet} f (x) d \nu_ {\alpha} (x).\tag{6}
\]

对 (5) 的第二成员有类似的关系，因此对 \(\alpha\) 求和便从 (5) 得出 (1)。

推论. --- X 的子集 N 为局部 \(\nu\) -可忽略的，必须且只须 \(\overline{\pi}^{-1}(N)\) 与点 \(t \in T\) （满足 \(g(t) > 0\) 者）所成的集的交是局部 \(\mu\) -可忽略的。

命题 2. --- 设 \(\pi\) 是 T 到 X 中的一个逆紧连续映射（GT, I, §10, No.~1），\(g\) 是 T 上的一个连续有限数值函数，使 \(g(t) > 0\) 对所有 \(t \in \mathrm{T}\) 成立。则对 \((\pi, g)\) 是 \(\mu\) -适应的，且若令

\[
\nu = \int g (t) \varepsilon_ {\pi (t)} d \mu (t),
\]

那么对每个数值函数 \(f \geqslant 0\) （定义在 \(X\) 上），

\[
\int^ {*} f (x) d \nu (x) = \int^ {*} f (\pi (t)) g (t) d \mu (t).\tag{7}
\]

显然 \(\pi\) 和 g 都是 \(\mu\) -可测的；此外，对每个函数 \(\psi \in \mathcal{K}(\mathrm{X})\) ，\(\psi \circ \pi\) 都是具有紧支集的连续函数，因为 \(\pi\) 是逆紧的；因此对 \((\pi, g)\) 是 \(\mu\) -适应的，而且映射 \(t \mapsto g(t) \varepsilon_{\pi(t)}\) 是淡连续的。

设 h 是 T 上的一个下半连续函数，使得

\[
f (\pi (t)) g (t) \leqslant h (t) \quad \text { for   all } t \in \mathrm{T}.
\]

我们将证明

\[
\int^ {*} f (x) d \nu (x) \leqslant \int^ {*} h (t) d \mu (t).\tag{8}
\]

由上积分的定义，这将蕴涵不等式

\[
\int^ {*} f (x) d \nu (x) \leqslant \int^ {*} f (\pi (t)) g (t) d \mu (t),
\]

把它与 §3, No.~2 的不等式 (7) 相结合，便证明了 (7)。

为证明 (8)，让我们定义一个函数 \(\overline{f}\) （\(X\) 上的）如下：\(\overline{f}(x)\) 是 \(h(t) / g(t)\) 在集 \(\overline{\pi}^{-1}(x)\) 中的下确界（下确界等于 \(+\infty\) ，如果 \(\overline{\pi}^{-1}(x) = \varnothing\) ）。函数 \(\overline{f}\) 具有下列性质：

\(1^{\circ}\overline{f} (x)\geqslant f(x)\) 对所有 \(x\in \mathrm{X}\) 成立（因为 \(g(t) > 0\) 对所有 \(t\in \mathrm{T}\) 成立）。

\[
2 ^ {\circ} \overline {{f}} (\pi (t)) g (t) \leqslant h (t) \text {   for   all   } t \in \mathrm{T}.
\]

\(3^{\circ}\) 函数 \(\overline{f}\) 是下半连续的，依据引理，因为函数 \(h / g\) 在 T 上是下半连续的。

因此，鉴于 §3, No.~1 的命题 2a)：

\[
\int^ {*} f (x) d \nu (x) \leqslant \int^ {*} \overline {{{f}}} (x) d \nu (x) = \int^ {*} \overline {{{f}}} (\pi (t)) g (t) d \mu (t) \leqslant \int^ {*} h (t) d \mu (t),
\]

这就建立了 (8)，并完成了证明。

\subsection{关于点测度积分的可测性}

命题 3. --- 设 \((\pi, g)\) 是一个 \(\mu\) -适应对，并令

\[
\nu = \int g (t) \varepsilon_ {\pi (t)} d \mu (t).
\]

设 \(f\) 是 \(X\) 到拓扑空间 \(G\) 中的一个映射，并设 \(S\) 是 \((\mu\text{-measurable})\) 集，由点 \(t \in T\) （满足 \(g(t) > 0\) 者）组成。\(f\) 为 \(\nu\) -可测的，必须且只须 \(f \circ \pi\) 在 \(S\) 上的限制是 \(\mu\) -可测的。

首先假设 \(f\) 是 \(\nu\) -可测的。由假设，集 \(\mathcal{K}\) （由紧子集 \(K\) （\(S\) 中的）组成，使 \(\pi\) 到 \(K\) 上的限制连续）是 \(\mu\) -稠密于 \(S\) 的（Ch. IV, §5, No.~10, 命题 15）。为证明 \(f \circ \pi\) 在 \(S\) 上的限制是 \(\mu\) -可测的，因此只需证明：对每个 \(K \in \mathcal{K}\) ，由紧子集 \(H\) （\(K\) 中的、使 \(f \circ \pi\) 在 \(H\) 上的限制连续者）组成的集是 \(\mu\) -稠密于 \(K\) 的（Ch. IV, §5, No.~8, 命题 13）。但由假设，存在紧集 \(\pi(K)\) 的一个分割，由一个 \(\nu\) -可忽略集 \(N\) 和一列紧集 \((C_n)\) 组成，使得 \(f\) 在每个 \(C_n\) 上的限制连续。在这些条件下，\(K \cap \overline{\pi}^{-1}(N)\) 与诸集 \(K \cap \overline{\pi}^{-1}(C_n)\) 构成 \(K\) 的一个分割；但 \(K \cap \overline{\pi}^{-1}(N)\) 依据 No.~2 的定理 1 的推论是 \(\mu\) -可忽略的，诸集 \(K \cap \overline{\pi}^{-1}(C_n)\) 是紧的，且 \(f \circ \pi\) 在后面这些集的每一个上的限制都连续，这就证明了 \(f \circ \pi\) 在 \(S\) 上的限制是 \(\mu\) -可测的。

反之，设情形如此；为证明 \(f\) 是 \(\nu\) -可测的，只需证明集 \(\mathfrak{L}\) （由紧子集 \(L\) （\(X\) 中的）组成，使 \(f\) 在 \(L\) 上的限制连续）是 \(\nu\) -稠密的（Ch. IV, §5, No.~10, 命题 15）。设 \(N\) 是 \(X\) 的一个子集，使 \(N \cap L\) 是 \(\nu\) -可忽略的（对每个 \(L \in \mathfrak{L}\) ），让我们证明 \(N\) 是局部 \(\nu\) -可忽略的。为此，我们必须证明 \(\overline{\pi}^{-1}(N) \cap S\) 是局部 \(\mu\) -可忽略的（No.~2 的定理 1 的推论）。现在，集 \(\mathfrak{H}\) （由紧子集 \(H\) （\(S\) 中的）组成，这些 \(H\) 使 \(\pi\) 和 \(f \circ \pi\) 在其上的限制都连续）由假设是 \(\mu\) -稠密于 \(S\) 的（Ch. IV, §5, No.~10, 命题 15）。因此只需证明 \(\overline{\pi}^{-1}(N) \cap H\) 是 \(\mu\) -可忽略的（对每个 \(H \in \mathfrak{H}\) ）。现在，\(\pi(H)\) 是紧的，并且可以与 \(H\) 按等价关系 \(\pi(t) = \pi(t')\) 所成的商空间等同，\(\pi\) 则等同于 \(H\) 到这个商空间上的典型映射（GT, I, §5, No.~2, 命题 3）。由于 \(f \circ \pi\) 在 \(H\) 上的限制连续，\(f\) 在 \(\pi(H)\) 上的限制就连续，换句话说 \(\pi(H) \in \mathfrak{L}\) ，从而 \(N \cap \pi(H)\) 是 \(\nu\) -可忽略的。由 No.~2 的定理 1 的推论，\(\overline{\pi}^{-1}(N \cap \pi(H)) \cap S\) 是局部 \(\mu\) -可忽略的；因此对集

\[
\mathrm{H} \cap \overline {{\pi}} ^ {- 1} (\mathrm{N}) \subset \overline {{\pi}} ^ {- 1} \bigl (\mathrm{N} \cap \pi (\mathrm{H}) \bigr) \cap \mathrm{S};
\]

但由于 H 是紧的，\(\mathrm{H} \cap \overline{\pi}^{-1}(\mathrm{N})\) 是 \(\mu\) -可忽略的，这就完成了证明。

注. --- 如果 f 是 X 到 Banach 空间 F 中的一个映射，那么说 \(f \circ \pi\) 在 S 上的限制是 \(\mu\) -可测的，与说函数 \((\mathbf{f} \circ \pi)g\) （定义在 T 上）是 \(\mu\) -可测的，是一回事，因为 \(g\) 是 \(\mu\) -可测的、在 S 中不为零且在 T - S 上为零（Ch. IV, §5, No.~10, 命题 15）。

\subsection{关于点测度积分的、取值于 Banach 空间函数的积分}

定理 2. --- 设 \((\pi, g)\) 是一个 \(\mu\) -适应对，并令

\[
\nu = \int g (t) \varepsilon_ {\pi (t)} d \mu (t).
\]

设 \(\mathbf{f}\) 是定义在 \(\mathbf{X}\) 上的一个函数，取值于 Banach 空间 \(\mathbf{F}\) 或 \(\overline{\mathbf{R}}\) 中。\(\mathbf{f}\) 为本质 \(\nu\) -可积的，必须且只须 \(t \mapsto \mathbf{f}(\pi(t))g(t)\) 是本质 \(\mu\) -可积的，此时

\[
\int \mathbf {f} (x) d \nu (x) = \int \mathbf {f} (\pi (t)) g (t) d \mu (t).\tag{9}
\]

此外假设 \(\pi\) 是连续且逆紧的，g 是连续的且使 \(g(t) > 0\) 对所有 \(t \in T\) 成立。那么，f 为 \(\nu\) -可积的，必须且只须 \(t \mapsto \mathbf{f}(\pi(t))g(t)\) 是 \(\mu\) -可积的。

\begin{enumerate}
\def\labelenumi{\Alph{enumi})}
\tightlist
\item
  我们首先处理测度 \(\mu\) 具有紧支集 K 且 g 在其上有界的情形。此时测度 \(\mu\) 和 \(\nu\) 都是有界的，因而可以把陈述中的“本质可积”换成“可积”。假设 f 是 \(\nu\) -可积的：则函数 \(\mathbf{f}(\pi(t))g(t)\) 是 \(\mu\) -可积的，且依据 §3, No.~3 的定理 1，关系式 (9) 成立。反之，假设 \(\mathbf{f}(\pi(t))g(t)\) 是 \(\mu\) -可积的：则 f 是 \(\nu\) -可测的（No.~3, 命题 3 及注），并且
\end{enumerate}

\[
\int^ {\bullet} | \mathbf {f} (x) | d \nu (x) = \int^ {\bullet} | \mathbf {f} (\pi (t)) | g (t) d \mu (t) <   + \infty
\]

（No.~2, 定理 1）；因此 \(\mathbf{f}\) 是本质 \(\nu\) -可积的（§1, No.~3, 命题 9），从而是 \(\nu\) -可积的。于是 §3, No.~3 的定理 1 蕴涵 (9)。

\begin{enumerate}
\def\labelenumi{\Alph{enumi})}
\setcounter{enumi}{1}
\tightlist
\item
  现在转到一般情形。设 \(\mathfrak{K}\) 是由紧子集 \(K\) （\(T\) 中的、使 \(g|K\) 连续者）组成的集：\(\mathfrak{K}\) 是 \(\mu\) -稠密的（Ch. IV, §5, No.~10, 命题 15），因此测度 \(\mu\) 是一个测度族 \((\mu_{\alpha})_{\alpha \in A}\) 之和，这些测度的支集都是 \(\mathfrak{K}\) 的元素（§2, No.~3, 命题 4）。对 \((g,\pi)\) 显然是 \(\mu_{\alpha}\) -适应的（对每个 \(\alpha \in A\) ），且测度 \(\nu\) 是测度族 \(\nu_{\alpha} = \int \varepsilon_{\pi(t)} g(t) d\mu_{\alpha}(t)\) 之和（§3, No.~1, 命题 1 的推论）。由于 A) 的论证可以应用于测度 \(\mu_{\alpha}, \nu_{\alpha}\) ，陈述的第一部分便由 §2, No.~2 的命题 3 得出。
\end{enumerate}

函数 \(\mathbf{f}\) （相应地，\(t\mapsto \mathbf{f}(\pi (t))g(t)\) ）为关于 \(\nu\) （相应地，关于 \(\mu\) ）可积的，必须且只须它是本质可积的，并且

\[
\int^ {*} | \mathbf {f} (x) | d \nu (x) <   + \infty \quad (\mathrm{resp.} \int^ {*} | \mathbf {f} (\pi (t)) | g (t) d \mu (t) <   + \infty).
\]

因此陈述的第二部分由第一部分及命题 2 得出。

注. --- 设 \((\pi, g)\) 是一个 \(\mu\) -适应对，\(\pi'\) 是 T 到 X 中的一个映射，\(g'\) 是定义在 T 上的一个有限数值函数 \(\geqslant 0\) ，使得 \(\pi'\) （相应地，\(g'\) ）等于 \(\pi\) （相应地，g）局部几乎处处（关于 \(\mu\) ）。则对 \((\pi', g')\) 是 \(\mu\) -适应的，测度 \(\lambda_t = g(t)\varepsilon_{\pi(t)}\) 与 \(\lambda_t' = g'(t)\varepsilon_{\pi'(t)}\) 局部几乎处处相等，且 \(\int g(t)\varepsilon_{\pi(t)} d\mu(t) = \int g'(t)\varepsilon_{\pi'(t)} d\mu(t)\) 。现在，如果 \(\pi'\) 和 \(g'\) 只是局部几乎处处（关于 \(\mu\) ）有定义，并且存在一个 \(\mu\) -适应对 \((\pi, g)\) 使得 \(\pi'\) （相应地，\(g'\) ）局部几乎处处等于 \(\pi\) （相应地，g），我们仍称对 \((\pi', g')\) 是 \(\mu\) -适应的，并置

\[
\int g ^ {\prime} (t) \varepsilon_ {\pi^ {\prime} (t)} d \mu (t) = \int g (t) \varepsilon_ {\pi (t)} d \mu (t)
\]

（参见 §3, No.~3 的注）。当 \(\pi\) 和 \(g\) 仅被假设为局部几乎处处有定义时，定理 1、定理 2 及命题 3 的陈述仍然有效。

\section{由数值密度定义的测度}

\subsection{局部可积函数}

命题 1. --- 设 g 是 T 中局部几乎处处（关于正测度 \(\mu\) ）有定义的一个函数，取值于 Banach 空间 F（相应地，\(\overline{R}\)）中。下列性质等价：

\begin{enumerate}
\def\labelenumi{\alph{enumi})}
\item
  对每个点 \(t \in \mathbb{T}\) ，存在一个邻域 \(\mathbb{V}\) （\(t\) 的），使得函数 \(\mathbf{g}\varphi_{\mathbb{V}}\) 是 \(\mu\) -可积的。
\item
  函数 \(\mathbf{g}\) 是 \(\mu\) -可测的，且对每个紧集 \(\mathrm{K} \subset \mathrm{T}\) ，\(\int^{*} |\mathbf{g}| \varphi_{\mathrm{K}} d\mu < +\infty\) 。
\item
  对每个数值函数 \(h \in \mathcal{K}(\mathrm{T})\) ，gh 都是 \(\mu\) -可积的。
\end{enumerate}

让我们证明 a) 蕴涵 b)；因为，由局部化原理（Ch. IV, §5, No.~2, 命题 4），函数 g 是可测的。另一方面，对每个 \(t \in \mathbf{K}\) ，由假设存在一个邻域 \(V_t\) （\(t\) 的，在 \(T\) 中），使 \(\mathbf{g}\varphi_{V_t}\) 可积；因此 \(K\) 可以被有限个邻域 \(V_i\) （ \(1 \leqslant i \leqslant n\) ）覆盖，使诸函数 \(\mathbf{g}\varphi_{V_i}\) 都可积。由于 \(|\mathbf{g}| \varphi_K \leqslant \sum_{i=1}^{n} |\mathbf{g}| \varphi_{V_i}\) ，就有 \(\int^* |\mathbf{g}| \varphi_K d\mu < +\infty\) 。

其次，b) 蕴涵 c)，因为此时 gh 可测，且若 L 是 h 的紧支集，则 \(|gh| \leqslant \|h\|\cdot|g|\varphi_{L}\) ，故由假设 \(\int^{*}|gh|d\mu < +\infty\) ；于是由可积性判别法（Ch. IV, §5, No.~6, 定理 5），gh 可积。

最后，c) 蕴涵 a)。的确，对每个 \(t \in T\) ，取 V 为 t 的一个紧邻域。存在 T 到 {[}0,1{]} 中的一个连续映射 h，它在 V 上等于 1 且具有紧支集（Ch. III, §1, No.~2, 引理 1）；由假设 gh 可积，因此 \(\mathbf{g}\varphi_{\mathrm{V}} = (\mathbf{g}h)\varphi_{\mathrm{V}}\) 也可积（Ch. IV, §5, No.~6, 定理 5 的推论 3）。

定义 1. --- 在 T 中局部几乎处处（关于正测度 \(\mu\) ）有定义、取值于 Banach 空间 F（相应地，\(\overline{R}\)）中的函数 g，如果满足命题 1 的条件 a), b), c)，就称为关于 \(\mu\) 局部可积的（或局部 \(\mu\) -可积的）。如果 \(\theta\) 是一个复测度，局部 \(\theta\) -几乎处处有定义的函数 g 称为局部 \(\theta\) -可积的，如果它关于正测度 \(|\theta|\) 局部可积。

如果 \(\mathbf{g}\) 是局部 \(\theta\) -可积的，那么每个局部几乎处处等于 \(\mathbf{g}\) 的函数都是局部可积的。显然，两个局部可积函数之和是局部可积的。取值于 \(\mathbf{F}\) 、处处有定义且关于 \(\theta\) 局部可积的函数构成一个向量空间，记作 \(\mathcal{L}_{\mathrm{loc}}^{1}(\mathrm{T},\theta ;\mathrm{F})\) ；当 \(\mathbf{F} = \mathbf{R}\) 或 \(\mathbf{C}\) 时，如无歧义，常省去对 \(\mathbf{F}\) 的提及。这个空间将总是（除明确说明相反者外）赋以由半范数 \(\mathbf{g}\mapsto \int |\mathbf{g}\varphi_{\mathbf{K}}|d|\theta |\) 所定义的拓扑，其中 \(\mathbf{K}\) 跑遍 \(\mathbf{T}\) 的紧子集所成的集。相联的 Hausdorff 空间，即 \(\mathcal{L}_{\mathrm{loc}}^{1}(\mathrm{T},\theta ;\mathrm{F})\) 关于局部几乎处处为零的映射所成的子空间 \(\mathcal{N}_{\mathrm{F}}^{\infty}\) 的商，记作 \(\mathrm{L}_{\mathrm{loc}}^{1}(\mathrm{T},\theta ;\mathrm{F})\) 。空间 \(\mathrm{L}_{\mathrm{loc}}^{1}(\mathrm{T},\theta ;\mathrm{F})\) 与 \(\mathrm{L}_{\mathrm{loc}}^{1}(\mathrm{T},|\theta |; \mathrm{F})\) 是相同的。

可以证明，刚才定义的拓扑向量空间是完备的（习题 31）。

在每个紧集上本性有界的每个可测函数 g 都是局部可积的。对每个使 \(1 \leqslant p \leqslant +\infty\) 的数 p，每个函数 \(g \in L_{F}^{p}\) 都是局部可积的；的确，对每个函数 \(h \in \mathcal{K}(T)\) ，h 属于 \(L^{q}\) （其中 q 是与 p 共轭的指数），因此 gh 可积（Ch. IV, §6, No.~4, 定理 2 的推论 4）。

设 \(F, G, H\) 是三个 Banach 空间，\((\mathbf{u}, \mathbf{v}) \mapsto \Phi(\mathbf{u}, \mathbf{v})\) 是 \(F \times G\) 到 \(H\) 中的一个连续双线性映射。如果 \(f\) 局部可积且取值于 F，且 \(g \in L_{G}^{\infty}\) ，那么 \(\Phi(\mathbf{f}, \mathbf{g})\) 是局部可积的（Ch. IV, §6, No.~4, 定理 2 的推论 1）。

\subsection{由数值密度定义的测度}

设 \(g\) 是一个正数值函数，它局部 \(\mu\) -几乎处处（在 \(T\) 中）有定义且局部 \(\mu\) -可积；此时由 \(t\) （使 \(g(t) = +\infty\) 者）组成的集是局部 \(\mu\) -可忽略的，因为 \(g\varphi_{K}\) 是 \(\mu\) -可积的（对每个紧集 \(K\) ）（Ch. IV, §2, No.~3, 命题 7）。现在设 \(g'\) 是一个局部可积的函数，它是正的且有限，并等于 \(g\) 局部 \(\mu\) -几乎处处；置 \(\lambda_t' = g'(t)\varepsilon_t\) 。映射 \(t \mapsto \lambda_t'\) （\(T\) 到 \(\mathcal{M}_+(T)\) 中的）是淡 \(\mu\) -可测的且标量本质可积（或者再说一次，对 \((I,g')\) ——其中 I 是 \(T\) 的恒等映射——是 \(\mu\) -适应的）；积分 \(\nu = \int \lambda_t' d\mu(t)\) 不依赖于特定的函数 \(g'\) （局部几乎处处等于 \(g\) 、用于定义测度 \(\lambda_t'\) 的那一个）。测度 \(\nu\) 由下列条件确定

\[
\int f (t) d \nu (t) = \int f (t) g (t) d \mu (t) \quad \mathrm{for} f \in \mathcal {K} (\mathrm{T}).\tag{1}
\]

现在，如果 \(\theta\) 是一个复测度，且 u 是一个复函数（或取值于 \(\overline{R}\) 中的函数），它局部 \(\theta\) -几乎处处有定义且关于 \(\theta\) 局部可积，那么可以写

\[
\begin{array}{l} {u = g _ {1} - g _ {2} + i (g _ {3} - g _ {4})} \\ {\theta = \mu_ {1} - \mu_ {2} + i (\mu_ {3} - \mu_ {4})} \end{array}\tag{2}
\]

其中 \(\mu_{1} = (\mathcal{R}\theta)^{+}\) ，\(\mu_{2} = (\mathcal{R}\theta)^{-}\) ，\(\mu_{3} = (\mathcal{I}\theta)^{+}\) ，\(\mu_{4} = (\mathcal{I}\theta)^{-}\) （Ch. III, §1, No.~5），而 \(g_{1}, g_{2}, g_{3}, g_{4}\) 有类似的含义；由于 \(|u|\) 是局部 \(|\theta|\) -可积的，诸正函数 \(g_{i}\) （ \(i = 1, 2, 3, 4\) ）中的每一个对每个正测度 \(\mu_{j}\) （ \(j = 1, 2, 3, 4\) ）都是局部可积的，于是映射

\[
f \mapsto \int f (t) u (t) d \theta (t)
\]

在 \(\mathcal{K}(\mathrm{T})\) 上是一个复测度。

定义 2. --- 设 \(\theta\) 是一个复测度，\(u\) 是一个复函数（或取值于 \(\overline{\mathbf{R}}\) 中的函数），它局部 \(\theta\) -几乎处处有定义且局部 \(\theta\) -可积。T 上的复测度 \(f \mapsto \int f u d\theta\) 称为测度 \(\theta\) 与函数 \(u\) 的乘积，或以 \(u\) 为密度（关于 \(\theta\) ）的测度，记作 \(u \cdot \theta\) 。

每个作为一个正测度 \(\mu\) 与一个局部 \(\mu\) -可积函数之乘积的复测度，都称为以 \(\mu\) 为基的测度。

关系式 \(\eta = u \cdot \theta\) 按惯例又写成

\[
d \eta (t) = u (t) d \theta (t).
\]

当 u 处处有定义且连续时，就重新得到 Ch. III, §1, No.~4 中给出的定义。显然，如果 \(u_{1}\) 和 \(u_{2}\) 都是局部 \(\theta\) -可积的，那么 \((u_{1}+u_{2})\cdot\theta=u_{1}\cdot\theta+u_{2}\cdot\theta\) 。类似地，如果 \(\theta_{1}\) 和 \(\theta_{2}\) 是 T 上两个测度，且 u 是一个关于 \(\theta_{1}\) 和 \(\theta_{2}\) 都局部可积的函数，那么 u 关于 \(\theta_{1}+\theta_{2}\) 局部可积，且有 \(u\cdot(\theta_{1}+\theta_{2})=u\cdot\theta_{1}+u\cdot\theta_{2}\) 。

此后我们将限于处处有定义的函数的情形；推广到局部几乎处处有定义的函数，这总是明显的，留给读者。

下面的命题在大多数情形允许把复测度的情形化归为正测度的情形：

命题 2. --- 设 \(\theta\) 是一个复测度，\(u\) 是一个局部 \(\theta\) -可积的复函数；则

\[
\left| u \cdot \theta \right| = \left| u \right| \cdot \left| \theta \right|.\tag{3}
\]

我们先给出一个辅助结果：

引理 1. --- 设 \(\theta\) 是一个复测度，\(f\) 是 \(\overline{\mathcal{L}}_{\mathbf{C}}^{1}(\mathrm{T},\theta)\) 的一个元素。则

\[
\langle | \theta |, | f | \rangle = \sup _ {c \in \mathcal {K} _ {1}} | \langle \theta , c f \rangle | = \sup _ {c \in \mathcal {B} _ {1}} | \langle \theta , c f \rangle |,\tag{4}
\]

其中 \(\mathcal{K}_{1}\) （相应地，\(B_{1}\) ）表示具有紧支集的连续（相应地，Borel）复函数 c 中使 \(|c| \leqslant 1\) 者所成的集。

我们先处理 \(f \in \mathcal{K}(\mathrm{T};\mathbf{C})\) 的情形。显然

\[
\sup _ {c \in \mathcal {K} _ {1}} | \langle \theta , c f \rangle | \leqslant \sup _ {c \in \mathcal {B} _ {1}} | \langle \theta , c f \rangle | \leqslant \langle | \theta |, | f | \rangle
\]

（Ch. IV, §4, No.~2, 命题 2）。另一方面，设 \(g\) 是 \(\mathcal{K}(\mathrm{T};\mathbf{C})\) 中使 \(|g| \leqslant |f|\) 的一个元素；\(g\) 是元素序列 \((g_n)\) （诸元素属于 \(\mathcal{K}(\mathrm{T};\mathbf{C})\) ）的一致极限，这些元素的支集都含在开集 \(\mathrm{U}\) （由 \(t\) 中使 \(f(t) \neq 0\) 者组成）中，并且显然可以假设 \(|g_n| \leqslant |f|\) （对每个 \(n\) ）。置 \(c_n(t) = g_n(t)/f(t)\) （对 \(t \in \mathrm{U}\) ），置 \(c_n(t) = 0\) （对 \(t \notin \mathrm{U}\) ）；则 \(c_n \in \mathcal{K}_1\) ，\(g = \lim_{n \to \infty} c_n f\) ，因此 \(|\langle \theta, g \rangle| = \lim_{n \to \infty} |\langle \theta, c_n f \rangle|\) ，最后

\[
\sup _ {| g | \leqslant | f |, g \in \mathscr {K} (\mathrm{T}; \mathbf {C})} | \langle \theta , g \rangle | \leqslant \sup _ {c \in \mathscr {K} _ {1}} | \langle \theta , c f \rangle |.
\]

注意到这个不等式的第一个端等于 \(\langle|\theta|,|f|\rangle\) （Ch. III, §1, No.~6, 公式 (12)），便得出结论。

其次，用 \(f\) 表示 \(\overline{\mathcal{L}}_{\mathbf{C}}^{1}(\theta)\) 的一个元素，我们来证明 (4) 仍然成立：只须验证这个关系式的三项连续地依赖于 \(f\) （对 \(\overline{\mathcal{L}}_{\mathbf{C}}^{1}(\theta)\) 的拓扑而言），因为它们在稠密子空间 \(\mathcal{K}(\mathrm{T};\mathbf{C})\) 上相合。这由下列不等式立即得出，其中 \(f\) 和 \(f'\) 表示 \(\overline{\mathcal{L}}_{\mathbf{C}}^{1}(\theta)\) 的元素：

\[
| \langle | \theta |, | f | \rangle - \langle | \theta |, | f ^ {\prime} | \rangle | \leqslant \langle | \theta |, | f - f ^ {\prime} | \rangle = \overline {{\mathrm{N}}} _ {1} (f - f ^ {\prime})
\]

\[
| \langle \theta , c f \rangle - \langle \theta , c f ^ {\prime} \rangle | \leqslant \langle | \theta |, | c | | f - f ^ {\prime} | \rangle \leqslant \overline {{{\mathrm{N}}}} _ {1} (f - f ^ {\prime})
\]

对所有 \(c \in B_{1}\) 成立。引理就此证毕。

转到命题 2 的证明，我们把引理应用于函数 uf，这里 h 属于 \(\mathcal{K}_{+}(\mathrm{T})\) 。这就得出：

\[
\langle | \theta |, | u h | \rangle = \sup _ {c \in \mathscr {K} _ {1}} | \langle \theta , c u h \rangle | = \sup _ {c \in \mathscr {K} _ {1}} | \langle u \cdot \theta , c h \rangle | = \langle | u \cdot \theta |, h \rangle .
\]

然而，第一个端也等于

\[
\langle | \theta |, | u | h \rangle = \langle | u | \cdot | \theta |, h \rangle .
\]

因此，两个测度 \(|u| \cdot |\theta|\) 与 \(|u \cdot \theta|\) 相等。

推论. --- 设 \(g_{1}\) 和 \(g_{2}\) 是两个局部 \(\mu\) -可积的数值函数；则

\[
\inf (g _ {1} \cdot \mu , g _ {2} \cdot \mu) = \inf (g _ {1}, g _ {2}) \cdot \mu ; \quad \sup (g _ {1} \cdot \mu , g _ {2} \cdot \mu) = \sup (g _ {1}, g _ {2}) \cdot \mu .
\]

特别地，若 \(g\) 是一个局部 \(\mu\) -可积的数值函数，则

\[
(g \cdot \mu) ^ {+} = g ^ {+} \cdot \mu ; (g \cdot \mu) ^ {-} = g ^ {-} \cdot \mu .
\]

这由命题 2 以及 Ch. II, §1, No.~1 的公式 (6) 立即得出。

\subsection{关于以密度定义的测度的积分}

在本小节的陈述中，g 表示一个正数值函数，它处处有定义且局部 \(\mu\) -可积；\(\theta\) 表示一个复测度，而 u 表示一个局部 \(\theta\) -可积的复函数。

No.~2 中的注记表明，§4 的结果适用于测度 \(\nu = g \cdot \mu = \int g(t)\varepsilon_t d\mu(t)\) （尽管测度 \(g(t)\varepsilon_t\) 仅当 \(g(t) \neq +\infty\) 时才有定义）。这样我们得到以下的陈述：

命题 3. --- 对每个数值函数 \(f \geqslant 0\) （定义在 T 上），

\[
\int^ {\bullet} f d \nu = \int^ {\bullet} (f g) d \mu .\tag{5}
\]

这由 §4, No.~2 的定理 1 得出。

推论 1. --- 为使一个在 Banach 空间中或在 \(\overline{R}\) 中取值的函数 f 对测度 \(u \cdot \theta\) 是局部可忽略的，必须且只须 uf 对 \(\theta\) 是局部可忽略的。

说 f （相应地，\(uf\) ）对 \(u \cdot \theta\) （相应地，对 \(\theta\) ）局部可忽略，等价于说 \(|f|\) （相应地，\(|u| \cdot |f|\) ）对 \(|u \cdot \theta|\) （相应地，对 \(|\theta|\) ）局部可忽略。于是，依据 No.~2 的命题 2，我们被化归到 \(f, u, \theta\) 均为正的情形；这时陈述由命题 3 立即得出。

推论 2. --- 设 \(u_1\) 和 \(u_2\) 是两个局部 \(\theta\) -可积的复函数。为使 \(u_1 \cdot \theta = u_2 \cdot \theta\) ，必须且只须 \(u_1\) 和 \(u_2\) 局部几乎处处相等。

我们立即被化归到证明 \(u \cdot \theta = 0\) 蕴涵 \(u(t) = 0\) 局部几乎处处成立；而 \(u \cdot \theta = 0\) 意味着函数 1 对测度 \(u \cdot \theta\) 局部可忽略。然后应用推论 1 即可。

推论 3. --- 设 u 是一个复函数，它对正测度 \(\mu\) 局部可积。为使 \(u \cdot \mu\) 是一个正测度，必须且只须 \(u(t) \geqslant 0\) 局部几乎处处成立。

因为，\(u \cdot \mu\) 是正的当且仅当 \(u \cdot \mu = |u \cdot \mu| = |u| \cdot \mu\) （命题 2），而这等价于 \(u = |u|\) 局部几乎处处成立（推论 2）。

命题 4. --- 为使一个映射 f （T 到拓扑空间 G 中的）对测度 \(u \cdot \theta\) 可测，必须且只须 f 到 \(\theta\) -可测集 S（由 t 中使 \(u(t) \neq 0\) 者组成）上的限制是 \(\theta\) -可测的。

当 u 和 \(\theta\) 为正时，这由 §4, No.~3 的命题 3 立即得出。然后，借助命题 2，结果推广到 u 和 \(\theta\) 为复的情形。

推论. --- 设 f 是定义在 T 上的一个函数，取值于 Banach 空间 F 或 \(\overline{R}\) 。为使 f 是 \((u\cdot\theta)\) -可测的，必须且只须 uf 是 \(\theta\) -可测的。

因为，\(u\mathbf{f}\) 是 \((uf)|S\) 到 \(T\) 上的用 0 所作的延拓。

定理 1. --- 设 f 是定义在 T 上的一个函数，取值于 Banach 空间 F 或 \(\overline{R}\) 。为使 f 对测度 \(\eta = u\cdot \theta\) 本质可积，必须且只须 uf 本质 \(\theta\) -可积，此时

\[
\int \mathbf {f} d \eta = \int (u \mathbf {f}) d \theta .\tag{6}
\]

再设 \(u\) 连续且 \(u(t) \neq 0\) 对所有 \(t \in \mathbb{T}\) 成立；则 \(\mathbf{f}\) 对测度 \(\eta\) 可积当且仅当 \(u\mathbf{f}\) 对 \(\theta\) 可积。

u 和 \(\theta\) 为正的情形由 §4, No.~4 的定理 2 立即得出。于是陈述的第一个和最后一个断论随之立即成立，因为 f 关于 \(\eta = u \cdot \theta\) 本质可积（相应地，可积）当且仅当它关于 \(|\eta| = |u| \cdot |\theta|\) 本质可积（相应地，可积）。最后，设 f 对 \(\eta\) （从而对 \(|\eta|\) ）本质可积；我们利用分解式 (2)：f 对每个测度 \(\eta_{ij} = g_i \cdot \mu_j\) （i = 1, 2, 3, 4, j = 1, 2, 3, 4）本质可积，因为这些测度都 \(\leqslant |\eta|\) 。我们有

\[
\int \mathbf {f} d \eta_ {i j} = \int g _ {i} \mathbf {f} d \mu_ {j}.
\]

公式 (6) 由此立即得出。

推论. --- 为使测度 \(u \cdot \theta\) 有界，必须且只须 u 本质 \(\theta\) -可积。

例子. --- 设 A 是 T 的一个子集；为使 \(\varphi_{A}\) 局部 \(\mu\) -可积，必须且只须 A 是 \(\mu\) -可测的。假设该条件成立，置 \(\nu = \varphi_{A} \cdot \mu\) ；则对每个数值函数 \(f \geqslant 0\) （定义在 T 上），我们有

\[
\int^ {\bullet} f d \nu = \int^ {\bullet} f \varphi_ {\mathrm{A}} d \mu ,
\]

这个量也记作 \(\int_{\mathrm{A}}^{\bullet}f d\mu\) （参见 Ch. IV, §5, No.~6）。为使一个映射 \(g\) （T 到拓扑空间 G 中的）是 \(\nu\) -可测的，必须且只须 \(g\) 在 A 上的限制是 \(\mu\) -可测的。为使一个映射 \(\mathbf{f}\) （T 到 Banach 空间 F 或 \(\overline{\mathbf{R}}\) 中的）本质 \(\nu\) -可积，必须且只须 \(\mathbf{f}\varphi_{\mathrm{A}}\) 本质 \(\mu\) -可积，此时

\[
\int \mathbf {f} d \nu = \int \mathbf {f} \varphi_ {\mathrm{A}} d \mu ,
\]

这个表达式也记作 \(\int_{A} f d\mu\) 。注意，若 T 到 G （相应地，F、\(\overline{R}\) ）中的两个映射在 A 上一致，则其中一个是 \(\nu\) -可测的（相应地，本质 \(\nu\) -可积的）必须且只须另一个也如此。现在若 g 是 T 的任意子集 \(B \supset A\) 到 G 中的一个映射，则称 g 在 A 上是 \(\mu\) -可测的，如果 g 在 A 上的限制到 T 上的任意延拓是 \(\nu\) -可测的；这等价于说 g 在 A 上的限制是 \(\mu\) -可测的。B 到 Banach 空间 F 或 \(\overline{R}\) 中的一个映射 f 称为在 A 上本质 \(\mu\) -可积的，如果 f 在 A 上的限制到 T 上的某个延拓 \(\overline{f}\) 是本质 \(\nu\) -可积的；此时置

\[
\int_ {\mathrm{A}} \mathbf {f} d \mu = \int_ {\mathrm{A}} \overline {{\mathbf {f}}} d \mu = \int \overline {{\mathbf {f}}} \varphi_ {\mathrm{A}} d \mu ,
\]

并称 \(\int_{\mathrm{A}} \mathbf{f} d\mu\) 为 \(\mathbf{f}\) 在 \(\mathrm{A}\) 上的积分（或延拓到 \(\mathrm{A}\) 的积分）。若 \(f\) 是一个数值函数 \(\geqslant 0\) （定义在 \(\mathrm{B} \supset \mathrm{A}\) 上），则类似地定义 \(\int_{\mathrm{A}}^{*} f d\mu\) 和 \(\int_{\mathrm{A}}^{\bullet} f d\mu\) 。最后，数值函数 \(g\) （定义在 \(\mathrm{B} \supset \mathrm{A}\) 上）称为局部 \(\mu\) -可积于 \(\mathrm{A}\) 的，如果 \(\overline{g}\) （到 \(\mathrm{T}\) 上的一个延拓，由 \(g\) 在 \(\mathrm{A}\) 上的限制扩充而成）是局部 \(\nu\) -可积的：这等价于说，对每个紧子集 \(\mathrm{K}\) （\(\mathrm{T}\) 中的），\(\overline{g}\varphi_{\mathrm{K} \cap \mathrm{A}}\) 是 \(\mu\) -可积的。

\subsection{乘积关于通常运算的行为}

命题 5. --- 设 \((\lambda_{\alpha})_{\alpha\in\mathbb{A}}\) 是 T 上的一个正测度族，关于关系 \(\leqslant\) 有向，在 \(\mathcal{M}(\mathrm{T})\) 中有一个上确界 \(\lambda\) 。为使正数值函数 g 局部 \(\lambda\) -可积，必须且只须 g 是局部 \(\lambda_{\alpha}\) -可积的（对每个 \(\alpha\in A\) ）且族 \((g\cdot\lambda_{\alpha})\) 在 \(\mathcal{M}(\mathrm{T})\) 中有上界；此时

\[
g \cdot \lambda = \sup _ {\alpha \in \mathrm{A}} g \cdot \lambda_ {\alpha}.
\]

条件的必要性是明显的。反之，设 \(g\) 对每个测度 \(\lambda_{\alpha}\) 都局部可积且族 \((g \cdot \lambda_{\alpha})_{\alpha \in \mathbb{A}}\) 有上界；用 \(\lambda'\) 表示它的上确界。函数 \(g\) 于是是 \(\lambda\) -可测的（§1, No.~4, 命题 11 的推论 2）；此外，对每个函数 \(h \in \mathcal{K}_{+}(\mathbf{T})\) ，

\[
\int^ {\bullet} (h g) d \lambda = \sup _ {\alpha \in A} \int^ {\bullet} (h g) d \lambda_ {\alpha} = \sup _ {\alpha \in A} \int^ {\bullet} h d (g \cdot \lambda_ {\alpha}) = \int^ {\bullet} h d \lambda^ {\prime}
\]

（§1, No.~4, 命题 11）。这首先蕴涵第一个端对任何 \(h\) 都是有限的，从而 \(g\) 局部 \(\lambda\) -可积；因此符号 \(\int^{\bullet}\) 可以换成 \(\int\) ，公式可以写成 \(\int h d(g \cdot \lambda) = \int h d\lambda'\) 。由此得出 \(g \cdot \lambda = \lambda'\) ，证明就此完成。

推论. --- 假设 \(\mu\) 是 T 上一个测度族 \((\mu_{\alpha})_{\alpha\in\mathrm{A}}\) 的和。为使定义在 T 上的正数值函数 g 局部 \(\mu\) -可积，必须且只须 g 是局部 \(\mu_{\alpha}\) -可积的（对每个 \(\alpha\in A\) ）且族 \((g\cdot\mu_{\alpha})_{\alpha\in\mathrm{A}}\) 可和。此时

\[
g \cdot \mu = \sum_ {\alpha \in A} g \cdot \mu_ {\alpha}.\tag{7}
\]

设 \((g_{\alpha})_{\alpha \in A}\) 是定义在 T 上的一个 \(\mu\) -可测正函数族。设 \(S_{\alpha}\) 是由 \(t \in T\) 中使 \(g_{\alpha}(t) \neq 0\) 者组成的集。我们称族 \((g_{\alpha})\) 是局部可数的，如果族 \((S_{\alpha})\) 是局部可数的（Ch. IV, §5, No.~9）；这等价于说，对 T 中每个紧集 K，由 \(\alpha \in A\) （使 \(g_{\alpha}|K\) 不为零者）组成的集是可数的。

命题 6. --- 设 \((g_{\alpha})_{\alpha \in A}\) 是定义在 T 上的局部 \(\mu\) -可积正数值函数的一个局部可数族。为使函数 \(g = \sum_{\alpha \in A} g_{\alpha}\) 局部 \(\mu\) -可积，必须且只须测度族 \((g_{\alpha} \cdot \mu)_{\alpha \in A}\) 可和，此时

\[
g \cdot \mu = \sum_ {\alpha \in A} g _ {\alpha} \cdot \mu .\tag{8}
\]

明显地，\(g\) 是 \(\mu\) -可测的（Ch. IV, §5, No.~2, 命题 4 及 No.~4, 定理 2 的推论 1）。因此，为使 \(g\) 局部 \(\mu\) -可积，必须且只须 \(\mu^{\bullet}(gf)\) 对每个 \(f \in \mathcal{K}_{+}(\mathrm{T})\) 都是有限的。现在，由于由 \(\alpha \in \mathrm{A}\) （使 \(g_{\alpha}f \neq 0\) 者）组成的集是可数的，我们有 \(\mu^{\bullet}(gf) = \sum_{\alpha \in \mathrm{A}} \mu^{\bullet}(g_{\alpha}f)\) （§1, No.~1, 命题 2 的推论）。置 \(\nu_{\alpha} = g_{\alpha} \cdot \mu\) ；条件 \(\mu^{\bullet}(gf) < +\infty\) 等价于条件 \(\sum_{\alpha \in \mathrm{A}} \nu_{\alpha}(f) < +\infty\) ：换言之，\(g\) 局部 \(\mu\) -可积当且仅当族 \((\nu_{\alpha})\) 可和。用 \(\nu\) 表示这个族的和，前面的计算给出等式 \(\nu(f) = \mu^{\bullet}(gf)\) ，它等价于 (8)。

推论. --- 设 \((g_{\alpha})\) 是局部 \(\mu\) -可积数值函数的一个序列，使得测度序列 \(g_{n} \cdot \mu\) 递增。为使这个序列在 T 上实测度所成的序向量空间 \(\mathcal{M}(T)\) 中有上界，必须且只须函数 \(g = \sup g_{n}\) 局部 \(\mu\) -可积；此时在 \(\mathcal{M}(T)\) 中序列 \((g_{n} \cdot \mu)\) 的上确界是测度 \(g \cdot \mu\) 。

只须把命题 6 应用于函数（局部几乎处处为正的）\(g_n' = g_{n+1} - g_n\) 。

命题 7. --- 设 X 是一个在无穷远处可数的局部紧空间，并设 \(t \mapsto \lambda_t\) 是一个 \(\mu\) -适宜映射（T 到 \(\mathcal{M}_+(X)\) 中的）。

设 \(g\) 是定义在 \(X\) 上的一个正数值函数，对测度 \(\nu = \int \lambda_t d\mu(t)\) 局部可积。于是由 \(t \in T\) 中使 \(g\) 不是局部 \(\lambda_t\) -可积者组成的集对 \(\mu\) 局部可忽略，映射 \(t \mapsto g \cdot \lambda_t\) （局部 \(\mu\) -几乎处处有定义）是 \(\mu\) -适宜的，且

\[
g \cdot \nu = \int (g \cdot \lambda_ {t}) d \mu (t).\tag{9}
\]

设 \((\mathrm{K}_{n})_{n\in\mathbf{N}}\) 是 X 的紧子集的一个递增序列，它们的内部覆盖 X；若 \(\eta\) 是 X 上任一正测度，则说 g 局部 \(\eta\) -可积等价于说 \(g\varphi_{\mathrm{K}_{n}}\) 对每个 n 都是 \(\eta\) -可积的。现在设 \(\mathrm{H}_{n}\) 是由 \(t\in\mathrm{T}\) 中使 \(g\varphi_{\mathrm{K}_{n}}\) 不是 \(\lambda_{t}\) -可积者组成的集，并设 \(\mathrm{H}=\bigcup_{n}\mathrm{H}_{n}\) ；由于 \(\mathrm{H}_{n}\) 对所有 n 都局部 \(\mu\) -可忽略（§3, No.~3, 定理 1），H 亦然，这就证明了陈述的第一个断言。把 \(\lambda_{t}\) 对 H 中的 t 换成 0（这不改变测度 \(\nu\) ），我们就可以假设 g 是局部 \(\lambda_{t}\) -可积的（对每个 \(t\in\mathrm{T}\) ）。对每个定义在 X 上的 \(\nu\) -可测正函数 h，我们依命题 3 及 §3, No.~2 的命题 5 有

\[
\int^ {\bullet} h d (g \cdot \nu) = \int^ {\bullet} (g h) d \nu = \int^ {\bullet} d \mu (t) \int^ {\bullet} (g h) d \lambda_ {t} = \int^ {\bullet} d \mu (t) \int^ {\bullet} h d (g \cdot \lambda_ {t}).
\]

这个公式与 §3, No.~2 的命题 5 首先表明（取 \(h \in \mathcal{K}_{+}(\mathbf{X})\) ）：映射 \(t \mapsto g \cdot \lambda_{t}\) 是标量本质 \(\mu\) -可积的，且其积分为 \(g \cdot \nu\) ；换言之，关系式 (9) 成立。其次，把 \(\mu\) 换成一个正测度 \(\mu' \leqslant \mu\) ，并取 \(h\) 为一个正的下半连续函数：由这些关系式立即得出 \(t \mapsto g \cdot \lambda_{t}\) 是 \(\mu\) -适宜的（§3, No.~1, 定义 1）。

命题 8. --- 设 \(\theta\) 是 T 上一个复测度，\(g_{1}\) 是一个局部 \(\theta\) -可积的复函数，\(\theta_{1}\) 是测度 \(g_{1} \cdot \theta\) 。为使定义在 T 上的复函数 \(g_{2}\) 局部 \(\theta_{1}\) -可积，必须且只须 \(g_{2} g_{1}\) 局部 \(\theta\) -可积，此时

\[
g _ {2} \cdot \theta_ {1} = g _ {2} \cdot (g _ {1} \cdot \theta) = (g _ {2} g _ {1}) \cdot \theta\tag{10}
\]

（`结合律公式'）。

依命题 4 的推论，说 \(g_{2}\) 是 \(\theta_{1}\) -可测的等价于说 \(g_{2}g_{1}\) 是 \(\theta\) -可测的。设这个条件成立。对每个函数 \(f \in \mathcal{K}_{+}(T)\) ，我们依命题 2 和命题 3 有

\[
\int^ {\bullet} | g _ {2} | f d | \theta_ {1} | = \int^ {\bullet} | g _ {2} | f | g _ {1} | d | \theta | = \int^ {\bullet} | g _ {2} g _ {1} | f d | \theta |.
\]

因此，说 \(g_{2}\) 局部 \(\theta_{1}\) -可积等价于说 \(g_{2}g_{1}\) 局部 \(\theta\) -可积。设这个条件成立，依定理 1 我们有

\[
\int f d \left(g _ {2} \cdot \theta_ {1}\right) = \int f g _ {2} d \theta_ {1} = \int f g _ {2} g _ {1} d \theta = \int f d \left(g _ {2} g _ {1} \cdot \theta\right),
\]

这个公式等价于 (10)。
